\documentclass[11pt]{article}

\usepackage[margin=1in]{geometry}
\usepackage[brazilian]{babel}
\usepackage{amsmath,amssymb,amsthm,mathtools}
\usepackage{enumitem}
\usepackage{hyperref}
\hypersetup{colorlinks=true,linkcolor=blue,citecolor=blue,urlcolor=blue}
\usepackage{orcidlink}
\usepackage{fontawesome5}

\newtheorem{theorem}{Teorema}[section]
\newtheorem{proposition}[theorem]{Proposição}
\newtheorem{lemma}[theorem]{Lema}
\newtheorem{corollary}[theorem]{Corolário}
\newtheorem{conjecture}[theorem]{Conjectura}
\theoremstyle{definition}
\newtheorem{definition}[theorem]{Definição}
\theoremstyle{remark}
\newtheorem{remark}[theorem]{Observação}
\newtheorem{example}[theorem]{Exemplo}

\DeclareMathOperator{\OpOrd}{ord}
\DeclareMathOperator{\Cov}{Cov}
\DeclareMathOperator{\Var}{Var}
\DeclareMathOperator{\Corr}{Corr}
\newcommand{\ZZ}{\mathbb{Z}}
\newcommand{\NN}{\mathbb{N}}
\newcommand{\RR}{\mathbb{R}}
\newcommand{\CC}{\mathbb{C}}
\newcommand{\Zp}{\mathbb{Z}_3}
\newcommand{\Zpu}{\mathbb{Z}_3^{\times}}
\newcommand{\e}[1]{e(#1)}

\title{A Generalização $qx+1$ da Medida de Syracuse de Tao:\\
Uma Equação de Pressão em Forma Fechada, uma Barreira de Endogenia,\\
e sua Caracterização Convergente por Sete Vias Independentes}

\author{Renato Augusto Tavares~{\small\href{mailto:dr.renatotavares@gmail.com}{\textcolor{black}{\faEnvelope}}}~\orcidlink{0009-0002-0196-3311}\\
\normalsize Universidade Federal de Goiás}
\date{\today}

\begin{document}

\maketitle

\begin{abstract}
Estudamos uma generalização natural para mapas $qx+1$ da variável
aleatória de Syracuse de Tao \cite{Tao2022}, o objeto probabilístico
que sustenta o resultado parcial incondicional mais forte já obtido em
direção à conjectura de Collatz. Provamos uma identidade exata, $\rho_{\mathrm{ann}}(\alpha) =
q^{\alpha-1}/(2^\alpha-1)$, para a pressão \emph{anelada} do funcional
populacional multiplicativo natural na árvore reversa do mapa $qx+1$
acelerado --- a taxa de crescimento exponencial, em média sobre todas
as raízes-residuais numa dada profundidade, da função de partição
associada --- via uma bijeção exata de contagem de fibra, corrigindo
um argumento de autômato finito para a mesma identidade que mostramos
não estar bem posto. A raiz não trivial de $\rho_{\mathrm{ann}}(\alpha)=1$
sofre uma transição estrutural em $q=5$: vale exatamente $\alpha^\ast=2$
em $q=3$ e cai estritamente abaixo de $1$ para todo $q\ge5$,
sinalizando uma mudança qualitativa de densidade natural positiva para
densidade nula da árvore reversa. Mostramos, além disso, que as duas
raízes dessa equação sempre ficam em lados opostos de uma transição de
congelamento quenched/anelado no sentido de polímeros dirigidos em
ambiente aleatório: a raiz menor está sempre não-congelada, então a
identidade anelada se transfere rigorosamente para a função de
contagem quenched real; a raiz maior está sempre congelada, então ela
não estabelece por si só o índice de cauda do fator de escala do
martingale que governa o termo de correção sub-exponencial, fato que
verificamos independentemente por cálculo exato. Em $q=3$ essa raiz
maior, $\alpha^\ast=2$, ainda coincide com o expoente de cauda
estabelecido empiricamente (independentemente do argumento de pressão)
por três métodos separados em trabalho anterior sobre o mapa de
Collatz clássico; para $q\ge5$ a alegação análoga fica como uma
conjectura explícita e fracamente sustentada. Em seguida
identificamos e caracterizamos formalmente uma \emph{barreira de
endogenia}, no sentido técnico de Aldous--Bandyopadhyay
\cite{AldousBandyopadhyay2005}, que obstrui qualquer prova apoiada
apenas na estrutura recursiva linear da árvore: soluções da forma
$G = W\cdot Y$, com $Y$ um fator limitado arbitrário invariante ao
longo da árvore, satisfazem a mesma recursão de ponto fixo que a
solução canônica (``endógena'') $W$ e são estatisticamente
indistinguíveis dela por qualquer observável que consigamos testar.
Mostramos, por seis vias tecnicamente não relacionadas, que fechar
essa lacuna equivale a uma única afirmação aritmética ainda em aberto:
a Weak Covering Conjecture de Wirsching de 1998 \cite{Wirsching1998},
a conjectura $\beta=1$ de Tao de 2020 \cite{Tao2020} (que mostramos ser
literalmente o mesmo objeto algébrico da Weak Covering Conjecture),
uma equação funcional para um análogo explícito em base $3$ da função
de Fabius que aparece no paper de Wirsching de 2003
\cite{Wirsching2003}, e dois lemas candidatos sugeridos como próximos
passos naturais pela técnica bivariada de decaimento de Fourier de Tao
\cite{Tao2022,Tao2011}, ambos executados até o fim e ambos falhando
pelo \emph{mesmo} ingrediente faltante: um produz um teorema de
estrutura genuíno em forma fechada junto com uma redução condicional
cuja hipótese (integrabilidade quadrática da densidade de Syracuse em
$\mathbb{Z}_3$) refutamos por um cálculo recursivo exato; o outro
produz uma dicotomia condicional genuína mas estritamente mais fraca,
junto com um resultado negativo exato, via identidade de Jensen,
mostrando que a maquinaria relevante de Littlewood--Offord/produtos de
Riesz tem uma folga intrínseca de orçamento de Fourier $\ell^1$ de um
fator $\approx 1{,}88$ no regime que importa. Complementamos esse
pacote teórico com um resultado puramente empírico e totalmente
independente: implementando a árvore reversa exata do mapa $5x+1$ e
aplicando extrapolação de Richardson a uma sequência de medições em
janelas encaixadas, resolvemos uma disputa antiga entre duas previsões
concorrentes para o expoente de contagem de $5x+1$
(Kontorovich--Lagarias \cite{KontorovichLagarias2009} contra um modelo
estocástico concorrente devido a Volkov, ambos discutidos em
\cite{KontorovichLagarias2009}), medindo $0{,}639$ (IC $95\%$
$[0{,}633,0{,}645]$) e excluindo decisivamente o valor concorrente
$0{,}678$. Por fim, uma busca literária dirigida situa a barreira com
precisão dentro da teoria mais ampla de decaimento de Fourier para
medidas autossimilares e autoconformais \cite{BakerBanaji2026,
Bremont}, e mostra que ela é estruturalmente inacessível a esse
ferramental por causa da rigidez dos subgrupos fechados de
$\mathbb{Z}_3^\times$. Argumentamos que essa convergência de sete
formulações independentes e tecnicamente não relacionadas sobre a
mesma afirmação não provada é, por si só, evidência significativa de
que a obstrução é uma característica genuína do problema, não um
artefato de um método isolado, e enunciamos com precisão o que resta
provar.
\end{abstract}

\tableofcontents

\section{Introdução}

\subsection{A conjectura de Collatz e a variável aleatória de Syracuse}

A conjectura de Collatz (ou $3x+1$) afirma que iterar o mapa
\[
  C(n) = \begin{cases} n/2 & n \text{ par} \\ 3n+1 & n \text{ ímpar} \end{cases}
\]
partindo de qualquer inteiro positivo $n$ eventualmente alcança $1$.
Apesar do enunciado elementar, a conjectura permanece aberta, e o
resultado incondicional mais forte em sua direção é devido a Tao
\cite{Tao2022}, que mostrou que \emph{quase todas} as órbitas (no
sentido de densidade logarítmica) atingem valores quase limitados. A
prova de Tao introduz a \emph{variável aleatória de Syracuse}: fixado
$n$, sejam $a_1,a_2,\dots,a_n$ variáveis aleatórias geométricas i.i.d.\
com $\Pr(a_i=k)=2^{-k}$ para $k\ge 1$, e defina
\begin{equation}\label{eq:syracuse-def}
  F_n(a) := \sum_{j=1}^n 3^{j-1}\, 2^{-a_{[1,j]}} \pmod{3^n}, \qquad
  a_{[1,j]} := a_1+\cdots+a_j,
\end{equation}
onde $2^{-1}$ denota o inverso multiplicativo de $2$ módulo $3^n$. A
variável aleatória $\mathrm{Syrac}(\mathbb{Z}/3^n\mathbb{Z}) :=
F_n(a_1,\dots,a_n)$ modela, em sentido preciso, o resíduo de uma
órbita de Collatz acelerada módulo $3^n$ após $n$ passos ímpares, e a
prova de Tao reduz o enunciado-alvo da conjectura a controlar a
distribuição dessa variável aleatória, em particular via uma
estimativa de decaimento de Fourier:

\begin{proposition}[Tao 2022, Prop.~1.17]\label{prop:tao117}
Para todo $A>0$ existe $C_A$ tal que para todo $n\ge 1$ e todo
$\xi\in\mathbb{Z}/3^n\mathbb{Z}$ com $3\nmid \xi$,
\[
  \bigl|\mathbb{E}\, e\bigl(\xi\, F_n(a)/3^n\bigr)\bigr| \le C_A\, n^{-A},
\]
com $C_A$ uniforme em $n$ e $\xi$ (sujeito apenas a $3\nmid\xi$).
\end{proposition}

Esta é uma estimativa de decaimento super-polinomial (embora não
exponencial), e é o núcleo técnico do teorema de Tao. Uma consequência
recursiva da construção \eqref{eq:syracuse-def}, que usamos
repetidamente adiante, é a seguinte identidade de \emph{autossimilaridade
exata}, verificada diretamente contra o texto de Tao:

\begin{proposition}[Tao 2022, eq.~(1.23)]\label{prop:tao123}
Para $k\le n$,
\[
  \mathrm{Syrac}(\mathbb{Z}/3^n\mathbb{Z}) \bmod 3^k \; \equiv\;
  \mathrm{Syrac}(\mathbb{Z}/3^k\mathbb{Z})
\]
como identidade em lei. Equivalentemente, escrevendo $F_n(a) =
2^{-a_1}\bigl(1+3\,F_{n-1}(a_2,\dots,a_n)\bigr) \bmod 3^n$, as $k$
coordenadas mais grosseiras de $\mathrm{Syrac}(\mathbb{Z}/3^n\mathbb{Z})$
dependem apenas de $(a_1,\dots,a_k)$, os dígitos \emph{adjacentes à
raiz} na indexação de \eqref{eq:syracuse-def}, não os mais profundos.
\end{proposition}

A convenção de orientação embutida na Proposição~\ref{prop:tao123}
(que $a_1$, não $a_n$, é a coordenada adjacente à raiz) é uma fonte
recorrente de erro ao traduzir a construção de Tao para a linguagem de
árvore reversa, e sinalizamos isso explicitamente ao longo do texto.

\subsection{Contribuições deste paper}

Construindo sobre \cite{Tao2022} e sobre o programa clássico de
densidade de predecessores de Wirsching \cite{Wirsching1998,
Wirsching2003}, apresentamos as seguintes contribuições.

\begin{enumerate}[label=(\roman*)]
\item \textbf{Uma equação de pressão em forma fechada para $qx+1$
  (\S\ref{sec:pressure}).} Mostramos que o expoente de cauda que
  controla o martingale populacional da árvore reversa do mapa $qx+1$
  acelerado satisfaz a equação exata $q^{\alpha-1}=2^\alpha-1$, com uma
  transição estrutural em $q=5$.
\item \textbf{Uma barreira de endogenia precisa (\S\ref{sec:endogeny}).}
  Exibimos uma família explícita de soluções \emph{não-endógenas} da
  recursão de ponto fixo da árvore, provamos que são estatisticamente
  indistinguíveis da canônica por pressão, índice de cauda, ou
  regressão linear contra a medida limite, e identificamos exatamente
  qual hipótese extra (quase-independência dos dígitos $3$-ádicos
  frescos entre subárvores irmãs) é necessária para descartá-las.
\item \textbf{Uma cota empírica sobre o acoplamento residual
  (\S\ref{sec:calibration}).} Um experimento de controle de três
  braços dá $\rho_{\mathrm{eff}} \lesssim 0{,}06$ (IC 95\%) para o
  acoplamento aritmético efetivo entre subárvores irmãs, convertendo a
  lacuna teórica de (ii) numa cota empírica quantitativa.
\item \textbf{Um resultado empírico incondicional (\S\ref{sec:kl}).}
  Resolvemos, usando enumeração exata e extrapolação de Richardson, uma
  disputa de quinze anos entre duas previsões concorrentes para o
  expoente de contagem de $5x+1$.
\item \textbf{Um teste computacional direto da Weak Covering
  Conjecture (\S\ref{sec:wcc}).}
\item \textbf{Três regimes de precisão convergentes para a extensão
  bivariada de Tao (\S\ref{sec:regimes}), e triangulação com duas
  outras formulações independentes da mesma obstrução
  (\S\ref{sec:triangulation}).}
\item \textbf{Dois lemas candidatos, executados até o fim
  (\S\ref{sec:lemmas}).} Ambos falham pelo mesmo ingrediente faltante,
  produzindo no caminho um teorema de estrutura genuíno, uma refutação
  computacional de uma condição suficiente natural, e um resultado
  negativo exato via identidade de Jensen.
\item \textbf{Uma contextualização na literatura (\S\ref{sec:context}).}
  Mostramos que a barreira é estruturalmente inacessível à teoria geral
  de decaimento de Fourier para medidas autossimilares, por uma razão
  aritmética precisa (rigidez dos subgrupos fechados de
  $\mathbb{Z}_3^\times$), e identificamos uma sétima formulação
  independente, puramente combinatória, da mesma obstrução em trabalho
  muito recente \cite{Chang2026}.
\end{enumerate}

Enquadramos este paper deliberadamente não como uma tentativa de prova
de qualquer parte da conjectura de Collatz, mas como uma caracterização
precisa de uma obstrução que recorre, essencialmente na mesma forma,
em toda técnica que nós ou a literatura existente trouxemos para esta
linha de ataque específica. Retomamos essa escolha de enquadramento na
\S\ref{sec:discussion}.

\section{Preliminares: o mapa $qx+1$ acelerado e sua árvore reversa}
\label{sec:setup}

Fixe um inteiro ímpar $q\ge 3$ coprimo com $2$. O \emph{mapa $qx+1$
acelerado} é
\[
  T_q(n) = \frac{qn+1}{2^{v_2(qn+1)}}, \qquad n \text{ ímpar},
\]
onde $v_2$ é a valuação $2$-ádica. Para $q=3$ este é o mapa de Collatz
acelerado padrão. A \emph{árvore reversa} enraizada num inteiro ímpar
$u$ coprimo com $q$ é construída adjungando, em cada nó ímpar $u$, um
filho $w_a := (2^a u - 1)/q$ para todo $a\ge 1$ tal que
$2^a u \equiv 1 \pmod q$; esse $w_a$ é automaticamente um inteiro ímpar
sempre que é inteiro, já que um numerador ímpar dividido exatamente por
um denominador ímpar dá um quociente ímpar. Escrevendo
$d:=\OpOrd_q(2)$ para a ordem multiplicativa de $2$ módulo $q$, a
admissibilidade de $a$ para um dado pai $u$ depende só de $u\bmod q$:
existe um único resíduo $a_0(u)\in\{1,\dots,d\}$ com
$2^{a_0(u)}u\equiv 1\pmod q$, e os expoentes admissíveis são exatamente
$a\equiv a_0(u)\pmod d$. Nós com $u\equiv 0\pmod q$ são
\emph{estéreis} (sem filhos admissíveis); para $q=3$ isso recupera o
fato conhecido de que múltiplos de $3$ têm subárvores reversas
triviais.

Nesta árvore estudamos o funcional populacional linear $G$ que
satisfaz a recursão de ponto fixo
\begin{equation}\label{eq:G-recursion}
  G(u) = \sum_{a \equiv a_0(u) \ (d)} q\cdot 2^{-a}\, G(w_a),
\end{equation}
que para $q=3$ é exatamente a recursão do tipo martingale que sustenta
a linha de investigação $G(v)/\mu$ resumida na \S\ref{sec:endogeny}
abaixo, e que é o análogo direto em árvore reversa de
\eqref{eq:syracuse-def}.

\section{A equação de pressão em forma fechada}
\label{sec:pressure}

\subsection{O operador de pressão: uma identidade anelada exata}

Uma primeira tentativa natural de organizar a recursão
\eqref{eq:G-recursion} numa operador de transferência é acompanhar só
a classe de resíduo do pai módulo $q$. Essa tentativa falha, e vale a
pena registrar exatamente por quê, já que uma versão anterior deste
paper afirmava algo assim que não resiste a escrutínio.

\begin{example}[O autômato ingênuo mod $q$ não está bem definido]
\label{ex:naive-fails}
Tome $q=3$ (logo $d=\OpOrd_3(2)=2$) e $a=2$, admissível para todo pai
$u\equiv1\pmod3$. Dois pais $u=1$ e $u=7$, ambos $\equiv1\pmod3$, dão
$w_2(1)=(4-1)/3=1$ e $w_2(7)=(28-1)/3=9$, com $w_2(1)\equiv1\pmod3$
mas $w_2(7)\equiv0\pmod3$: a classe de resíduo do filho depende de
$u\bmod9$, não só de $u\bmod3$. Em geral, escrevendo $u=u_0+qk$ com
$u_0=u\bmod q$, obtém-se $w_a\bmod q = [(2^au_0-1)/q\bmod q] +
[2^ak\bmod q]$: dependência explícita de $k=\lfloor u/q\rfloor\bmod q$.
Nenhum módulo finito $q^k$ conserta isso: o mesmo cálculo mostra que
$w_a\bmod q^k$ depende de $u\bmod q^{k+1}$, um dígito além do que o
conhecimento de $u$ mod $q^k$ fornece. Não existe, portanto, um
operador de transferência de estado finito bem definido rastreando uma
única geração da recursão em classes de resíduo de qualquer módulo
finito fixo.
\end{example}

O que \emph{está} bem definido, e é o substituto correto do quadro
ingênuo do Exemplo~\ref{ex:naive-fails}, é uma bijeção entre uma fibra
admissível inteira num nível da torre de dígitos e o anel de resíduos
completo um nível abaixo.

\begin{lemma}[Bijeção de fibra]\label{lem:fibre-bijection}
Fixe $a\ge1$ e $j\ge1$. O mapa $v\mapsto w_a(v):=(2^av-1)/q$ é uma
bijeção de $\{v\bmod q^j : 2^av\equiv1\!\!\pmod q\}$ (uma progressão
aritmética de $q^{j-1}$ resíduos mod $q^j$) sobre
$\mathbb{Z}/q^{j-1}\mathbb{Z}$.
\end{lemma}
\begin{proof}
Bem definido: $v\bmod q^j$ determina $2^av-1\bmod q^j$, logo
$w_a(v)\bmod q^{j-1}$. Injetivo: se $v\not\equiv v'\pmod{q^j}$ dentro
da progressão, escreva $v-v'=qt$ com $t\not\equiv0\pmod{q^{j-1}}$;
então $w_a(v)-w_a(v')=2^at\not\equiv0\pmod{q^{j-1}}$ já que $2$ é
unidade mod $q^{j-1}$. A sobrejetividade segue por ser uma injeção
entre dois conjuntos finitos de mesmo tamanho $q^{j-1}$.
\end{proof}

Iterando o Lema~\ref{lem:fibre-bijection} ao longo de um caminho de
profundidade $k$ $(a_1,\dots,a_k)$ recupera-se, em forma truncada,
exatamente a identidade autossimilar \eqref{eq:X-selfsimilar} da
\S\ref{sec:context} abaixo: toda sequência
$(a_1,\dots,a_k)\in(\mathbb{Z}_{\ge1})^k$, sem restrição de
admissibilidade sobre os $a_i$ individuais, é admissível para
\emph{exatamente uma} raiz-residual $u_0\bmod q^k$, a saber
$u_0\equiv\sum_{i=1}^k q^{i-1}2^{-S_i}\pmod{q^k}$,
$S_i:=a_1+\cdots+a_i$. Defina a função de partição em profundidade $k$
enraizada em $u_0$,
\[
  Z_k(\alpha;u_0) := \sum_{\substack{(a_1,\dots,a_k)\text{ admissível}\\ \text{a partir de } u_0}} \ \prod_{i=1}^k \bigl(q\cdot2^{-a_i}\bigr)^\alpha .
\]

\begin{theorem}[Identidade de pressão anelada exata]\label{thm:pressure}
Para todo $q\ge3$ ímpar, todo $k\ge1$, e todo $\alpha>0$,
\begin{equation}\label{eq:annealed-identity}
  \sum_{u_0\in\mathbb{Z}/q^k\mathbb{Z}} Z_k(\alpha;u_0) \;=\; \Bigl(\frac{q^\alpha}{2^\alpha-1}\Bigr)^{k},
\end{equation}
logo a \emph{média} de $Z_k(\alpha;\cdot)$ sobre todas as $q^k$
raízes-residuais vale exatamente $\bigl(q^{\alpha-1}/(2^\alpha-1)\bigr)^k$;
a pressão anelada associada é
\begin{equation}\label{eq:pressure-closed-form}
  \rho_{\mathrm{ann}}(\alpha) = \frac{q^{\alpha-1}}{2^\alpha-1},
\end{equation}
dando a equação de pressão em $\rho_{\mathrm{ann}}=1$:
\begin{equation}\label{eq:pressure-eq}
  q^{\alpha-1} = 2^{\alpha}-1.
\end{equation}
\end{theorem}
\begin{proof}
Pelo Lema~\ref{lem:fibre-bijection} (iterado), as sequências
admissíveis de profundidade $k$ enraizadas nos $q^k$ resíduos distintos
particionam o conjunto completo e irrestrito $(\mathbb{Z}_{\ge1})^k$,
sem sobreposição nem omissão. Somar $\prod_i(q2^{-a_i})^\alpha$ sobre
\emph{todas} as sequências em $(\mathbb{Z}_{\ge1})^k$ portanto se
fatora como $\bigl(\sum_{a=1}^\infty (q2^{-a})^\alpha\bigr)^k =
\bigl(q^\alpha/(2^\alpha-1)\bigr)^k$, que é exatamente o lado esquerdo
de \eqref{eq:annealed-identity} somado sobre a partição.
\end{proof}

Verificamos \eqref{eq:annealed-identity} independentemente por
enumeração direta para $q\in\{3,5,7,9\}$, $k\in\{1,2,3\}$, e
$\alpha\in\{0{,}26,0{,}5,1,1{,}37,2\}$, batendo a menos do erro de
truncamento da cauda (rapidamente convergente) da soma sobre $a$.

\begin{remark}[Conservação de massa como corolário]
O lado direito de \eqref{eq:annealed-identity} vale exatamente $1^k=1$
em $\alpha=1$ para todo $q$ (conservação de massa do processo de
ramificação). Logo $\alpha=1$ é sempre raiz de \eqref{eq:pressure-eq}.
\end{remark}

\subsection{Quenched vs.\ anelado: uma transição de congelamento}
\label{sec:freezing}

O Teorema~\ref{thm:pressure} é uma afirmação sobre a \emph{média} de
$Z_k(\alpha;\cdot)$ sobre raízes-residuais; não é, por si só, uma
afirmação sobre $Z_k(\alpha;u_0)^{1/k}$ para uma raiz $u_0$ única e
fixa --- que é o que de fato governa o martingale populacional $G(v)$
para um inteiro específico $v$. Se as duas coincidem é exatamente a
questão quenched-versus-anelado familiar de polímeros dirigidos em
ambiente aleatório e do modelo de energia aleatória de Derrida,
governada por um critério de congelamento padrão. Escrevendo
$P(\alpha):=\log\rho_{\mathrm{ann}}(\alpha)$ para a log-pressão anelada
e $s(\alpha):=P(\alpha)-\alpha P'(\alpha)$ para a entropia associada, as
taxas quenched e anelada coincidem onde $s(\alpha)>0$ (uma fase
\emph{não-congelada}) e divergem onde $s(\alpha)<0$ (uma fase
\emph{congelada}, na qual a taxa quenched é dominada por caminhos raros
de baixa entropia, não pelo grosso do ensemble).

\begin{proposition}[A raiz maior está sempre congelada]\label{prop:always-frozen}
Para todo $q\ge3$ ímpar, escrevendo $\alpha_-<\alpha_+$ para as duas
raízes positivas de $\rho_{\mathrm{ann}}(\alpha)=1$: $s(\alpha_-)>0$
(não-congelada) e $s(\alpha_+)<0$ (congelada), incondicionalmente.
\end{proposition}
\begin{proof}
Em qualquer raiz, $P(\alpha)=0$, logo $s(\alpha)=-\alpha P'(\alpha)$.
Como $\rho_{\mathrm{ann}}=e^P$ é log-convexa (\S\ref{sec:endogeny-addendum}),
$P$ em si é estritamente convexa; com exatamente dois zeros
$\alpha_-<\alpha_+$, necessariamente $P'(\alpha_-)<0<P'(\alpha_+)$ (a
função precisa estar decrescendo no primeiro zero e crescendo no
segundo, em lados opostos de seu mínimo único). Logo
$s(\alpha_-)=-\alpha_-P'(\alpha_-)>0$ e
$s(\alpha_+)=-\alpha_+P'(\alpha_+)<0$.
\end{proof}

Para $q\ge5$, onde $\alpha_+=1$ sempre, isso tem forma fechada exata:
em $\alpha=1$, $2^\alpha-1=1$ logo $\log(2^\alpha-1)=0$, e
$s(1)=2\log2-\log q=\log(4/q)$, negativo exatamente quando $q>4$.
Então o congelamento da raiz trivial para todo $q\ge5$ ímpar não é
coincidência numérica, é a mesma condição de sinal da observação sobre
o drift médio abaixo ($\log q-2\log2>0$): o drift que faz órbitas
expandirem em média é, a menos de um sinal, exatamente o déficit de
entropia que congela $\alpha_+=1$.

Calculamos o limiar de congelamento $\alpha_c(q)$ (a raiz de
$s(\alpha)=0$, que pela Proposição~\ref{prop:always-frozen} sempre fica
estritamente entre $\alpha_-(q)$ e $\alpha_+(q)$) diretamente:
\begin{center}
\begin{tabular}{c|c}
$q$ & $\alpha_c(q)$ \\ \hline
$3$ & $1{,}355$ \\
$5$ & $0{,}794$ \\
$7$ & $0{,}564$ \\
$9$ & $0{,}437$
\end{tabular}
\end{center}
confirmando $\alpha_-(q)<\alpha_c(q)<\alpha_+(q)$ em todos os casos,
como a Proposição~\ref{prop:always-frozen} exige. Também verificamos
numericamente, com uma recursão memoizada independente calculando
$Z_k(\alpha;u_0)$ exatamente numa raiz fixa (não a média), que o
congelamento não é só um cálculo formal de entropia, mas uma divergência
genuína entre crescimento quenched e anelado: em $q=3$, $u_0=1$ --- um
ciclo genuíno de período um da dinâmica de resíduos, já que $w_2(1)=1$
exatamente --- $Z_k(2;1)^{1/k}$ mostra oscilação com razão de média
geométrica bem abaixo da taxa anelada $\rho_{\mathrm{ann}}(2)=1$; em
$q=5$, $Z_k(1;u_0)^{1/k}$ em várias raízes $u_0\in\{1,2,3,4\}$
igualmente falha em acompanhar $\rho_{\mathrm{ann}}(1)=1$. Relatamos
isso como confirmação qualitativa de congelamento (quenched
estritamente abaixo do anelado), não como um ajuste quantitativo a
nenhuma assintótica específica de fase congelada, que não derivamos.

A consequência, pela Proposição~\ref{prop:always-frozen}, é
incondicional: o Teorema~\ref{thm:pressure}, embora exato, se
transfere rigorosamente só para a metade de \emph{expoente de
contagem} do Teorema~\ref{thm:transition} abaixo (a raiz menor
$\alpha_-(q)$, sempre não-congelada); ele nunca estabelece, para
nenhum $q$, o índice de cauda do fator de escala do martingale $W_u$
também alegado ali, já que a raiz maior $\alpha_+(q)$ que governaria
essa alegação está sempre congelada. Tratamos isso honestamente como a
Conjectura~\ref{conj:tail-index} abaixo, não como mais uma consequência
do Teorema~\ref{thm:pressure}.

\subsection{A transição estrutural em $q=5$}

Além da raiz trivial $\alpha=1$, a equação~\eqref{eq:pressure-eq} tem
exatamente mais uma raiz positiva $\alpha^\ast(q)$ (pela convexidade
logarítmica estrita de $\alpha\mapsto q^{\alpha-1}/(2^\alpha-1)$ em
$\alpha>0$, discutida na \S\ref{sec:endogeny-addendum} abaixo, a
equação $\rho(M_q(\alpha))=1$ tem no máximo duas raízes positivas).
Resolvemos $\alpha^\ast(q)$ independentemente por bisseção:

\begin{center}
\begin{tabular}{c|c|c}
$q$ & segunda raiz $\alpha^\ast(q)$ & interpretação \\ \hline
$3$ & $2{,}000000$ & exata \\
$5$ & $0{,}650919$ & \\
$7$ & $0{,}373501$ & \\
$9$ & $0{,}258108$ &
\end{tabular}
\end{center}

Em $q=3$ a segunda raiz é exatamente $\alpha^\ast=2$, coincidindo com
o expoente de cauda universal da árvore reversa clássica de Collatz
estabelecido por métodos independentes (estimador de Hill sobre
amostras de árvores reais, e teoria de valores extremos sobre máximos
de blocos) em etapas anteriores deste projeto. Para $q\ge 5$, porém, a
segunda raiz cai \emph{abaixo} de $1$ em vez de acima: uma transição
estrutural genuína, não uma perturbação numericamente pequena.
Escrevendo $\alpha_-<\alpha_+$ para as duas raízes, temos
$\{\alpha_-,\alpha_+\}=\{1,2\}$ em $q=3$ mas
$\{\alpha_-,\alpha_+\}=\{\alpha^\ast(q),1\}$ para $q\ge5$: a raiz
trivial passa a ser a \emph{maior} das duas.

\begin{theorem}[Transição estrutural: densidade]\label{thm:transition}
Para $q=3$ a árvore reversa de $T_q$ tem densidade natural positiva
(função de contagem $N_u(x)\sim x$); para todo $q\ge 5$ ímpar ela tem
densidade natural nula, com $N_u(x)\sim W_u\cdot x^{\alpha_-(q)}$ para
algum fator de escala aleatório $W_u$, $\alpha_-(q)<1$.
\end{theorem}

Pela Proposição~\ref{prop:always-frozen}, $\alpha_-(q)$ é sempre
não-congelada, então a identidade anelada do Teorema~\ref{thm:pressure}
se transfere para a função de contagem quenched real, nesta raiz,
consistente com a confirmação empírica direta abaixo e com a medição
estatisticamente calibrada da \S\ref{sec:kl}.

\begin{remark}[No que essa "transferência" se apoia]\label{rem:transfer-basis}
Somos precisos sobre o status dessa transferência, já que é exatamente
o tipo de passo que causou o erro corrigido nesta seção. Para o modelo
genuíno i.i.d.\ de passeio aleatório ramificado que a identidade
anelada do Teorema~\ref{thm:pressure} descreve, quenched=anelado em
toda a região não-congelada é um resultado clássico e rigoroso
\cite{Biggins1992}. Para nossa recursão aritmética, nenhum teorema
análogo é conhecido por nós: a lacuna é a mesma já identificada no
Teorema~\ref{thm:barrier} (quase-independência de dígitos frescos entre
subárvores irmãs, aqui numa forma quantitativa de grandes desvios em
vez de qualitativa), e o resultado aritmético genuíno mais próximo na
literatura é uma cota unilateral de densidade de
Applegate--Lagarias \cite{ApplegateLagarias1995I,ApplegateLagarias1995II},
depois refinada por Krasikov--Lagarias \cite{KrasikovLagarias2003} (já
citado na \S\ref{sec:triangulation}), não uma afirmação quenched
bilateral correspondente. Tratamos, portanto, a transferência, em toda
raiz onde a invocamos, como apoiada por (i) o teorema clássico do
modelo estocástico, (ii) confirmação numérica direta, e não por uma
prova para a própria árvore aritmética --- nomeando a lacuna com
precisão em vez de deixar o rigor do teorema do modelo estocástico
vazar por associação para a alegação aritmética.
\end{remark}

\begin{conjecture}[Índice de cauda do fator de escala]\label{conj:tail-index}
$W_u$ tem cauda regularmente variável de índice
$\alpha_+(q)/\alpha_-(q)$, onde $\alpha_+(q)$ é a raiz maior de
\eqref{eq:pressure-eq} ($\alpha_+(3)=2$; $\alpha_+(q)=1$ para todo
$q\ge5$).
\end{conjecture}

Este é exatamente o expoente previsto pela teoria clássica de pontos
fixos da transformação de suavização (multitipo)
\cite{Guivarch1990,Liu2000,Mentemeier2016} para o modelo estocástico
genuíno: a equação de índice de cauda $\mathbb{E}[\sum_i A_i^s]=1$
vira, para nossos pesos, $q^{\alpha_-s-1}=2^{\alpha_- s}-1$, resolvida
exatamente por $s=\alpha_+/\alpha_-$. Isso é razão adicional para levar
a conjectura a sério como previsão, ainda que, pela
Observação~\ref{rem:transfer-basis}, a teoria só cubra rigorosamente o
modelo estocástico, não nossa recursão aritmética.

Enunciamos isso como conjectura, não como parte do
Teorema~\ref{thm:transition}, precisamente porque a
Proposição~\ref{prop:always-frozen} mostra que $\alpha_+(q)$ está
\emph{sempre} congelada: a identidade de pressão anelada não a
estabelece para nenhum $q$, e ela precisa se apoiar em evidência
externa ao Teorema~\ref{thm:pressure}. Em $q=3$ essa evidência é real e
tripla: o estimador de Hill sobre amostras de árvores reais e a teoria
de valores extremos sobre máximos de blocos (ambos de etapas anteriores
deste projeto) dão $\alpha\approx2$ independentemente do argumento de
pressão, e a confirmação do expoente de contagem abaixo soma uma
terceira linha relacionada de apoio. Para $q\ge5$ a evidência
correspondente é marcadamente mais fraca: uma medição inicial via
estimador de Hill da cauda de $W_u$ para $q=5$ (600 raízes,
concordância a duas casas decimais com o previsto $1{,}536\ldots$) foi
depois mostrada, ao longo deste projeto, ser estatisticamente
não-confirmatória nesse tamanho de amostra (o erro-padrão real do
estimador é $\approx0{,}45$, uma ordem de magnitude maior que a
concordância alegada). Um teste subsequente, substancialmente maior
(5000 raízes, quatro níveis de headroom até $10^8$, intervalos de
confiança via bootstrap, e uma regressão rank-size independente;
\texttt{collatz-endogeny/sec3-pressure-equation}) é encorajador mas
ainda não decisivo: a estimativa é notavelmente estável entre níveis
de headroom, e na fração de cauda julgada mais equilibrada (5\%) cai
em $1{,}58$ (IC 95\% $[1{,}41,1{,}80]$), próximo do previsto
$1{,}536$ --- mas, como é típico de estimadores tipo Hill, o valor
pontual é instável conforme a escolha da fração de cauda (de $1{,}39$
em 10\% a $2{,}10$ em 1\%), de modo que não tratamos isso como
confirmação. Não temos conhecimento de nenhum teste empírico
totalmente decisivo da Conjectura~\ref{conj:tail-index} para $q\ge5$
(por exemplo, uma comparação pré-registrada tipo
Clauset--Shalizi--Newman contra modelos alternativos de cauda), e a
listamos como problema em aberto explícito na \S\ref{sec:conclusion}.

O drift médio $\log q - 2\log 2$ muda de sinal exatamente no limiar
$q=5$ que separa os dois casos do Teorema~\ref{thm:transition}:
negativo para $q=3$ (órbitas contraem em média), positivo para $q\ge5$
(órbitas expandem em média). Isso não é coincidência ao lado do cálculo
de congelamento da \S\ref{sec:freezing}: em $\alpha=1$ a entropia é
exatamente $s(1)=\log(4/q)$, então o mesmo sinal que decide contração
versus expansão de órbitas também decide se a raiz trivial $\alpha=1$
está congelada.

\subsection{Confirmação empírica em árvores reversas reais}

Enumeração independente por busca em profundidade das árvores reversas
verdadeiras (não estendidas) para $q=5$ e $q=7$, seguindo a regra de
admissibilidade exata da \S\ref{sec:setup}, confirma o
Teorema~\ref{thm:transition} diretamente: inclinações de contagem por
década de $3$ raízes amostradas independentemente deram
$\approx 0{,}62$--$0{,}66$ para $q=5$ (previsto $0{,}650919$) e
$\approx 0{,}36$--$0{,}39$ para $q=7$ (previsto $0{,}373501$). Uma
medição mais cuidadosa e estatisticamente calibrada para $q=5$ é o
assunto da \S\ref{sec:kl}.

\begin{remark}[Calibração bibliográfica]\label{rem:novelty-109}
Uma checagem de novidade dirigida contra Kontorovich--Lagarias
\cite{KontorovichLagarias2009}, Wirsching \cite{Wirsching1998}, e um
resultado rigoroso posterior na direção forward que exclui $q\ge5$
para a generalização análoga $p=2$
\cite{GoncalvesGreenfeldMadrid2022} estabeleceu que nem o valor
numérico $\alpha^\ast(5)=0{,}650919$ nem a transição qualitativa em
$q\ge 5$ são novos: Wirsching já havia conjecturado a transição
heuristicamente em 1998, e Kontorovich--Lagarias já haviam calculado a
mesma quantidade (denotada $\eta_{5,BP}$ no Teorema~8.10 deles) via
métodos de grandes desvios para um passeio aleatório ramificado, uma
rota diferente da nossa derivação via equação de pressão, mas
numericamente idêntica a ela. Dado a \S\ref{sec:freezing}, essa
diferença de rota pode importar mais do que uma mera derivação
alternativa: um argumento genuíno de grandes desvios para um passeio
aleatório ramificado é, no sentido apropriado, uma afirmação quenched,
enquanto nossa identidade de pressão (Teorema~\ref{thm:pressure}) é
anelada. Que as duas concordem numericamente em $\alpha_-(5)$ é
consistente com $\alpha_-(5)$ ser não-congelada
(Proposição~\ref{prop:always-frozen}), onde quenched e anelado devem
coincidir; isso não certifica por si só que nossa rota anelada
concordaria com uma rota quenched numa raiz congelada. O que sobrevive
como novo aqui é a própria equação em forma fechada~\eqref{eq:pressure-eq},
unificando a família para $q$ arbitrário numa única identidade
algébrica, junto com sua derivação anelada rigorosa e confirmação
numérica independente.
\end{remark}

\section{A barreira de endogenia}
\label{sec:endogeny}

\subsection{Liberdade de gauge na recursão da árvore}

A recursão \eqref{eq:G-recursion} é linear em $G$, e a linearidade tem
uma consequência estrutural imediata fácil de passar despercebida: se
$W$ denota a solução ``aritmética'' --- aquela de fato realizada por,
e mensurável em relação a, os dígitos $3$-ádicos da árvore --- e $Y$ é
\emph{qualquer} função positiva limitada invariante ao longo da árvore
(em particular, qualquer função de uma coordenada auxiliar carregada
sem mudança de pai para filho), então $X := W\cdot Y$ satisfaz
\emph{exatamente a mesma recursão}.

\begin{proposition}[Liberdade de gauge]\label{prop:gauge}
Seja $\mathbb{Z}_3^\ast\times\mathcal{Y}$ o espaço produto no qual a
dinâmica dos dígitos $3$-ádicos age no primeiro fator exatamente como
na árvore reversa, enquanto passa uma coordenada auxiliar
$y\in\mathcal{Y}$ sem mudança para todo filho. Para $Y:\mathcal{Y}\to
(0,\infty)$ limitada e não constante, a função $X(r,y):=W(r)\cdot Y(y)$
satisfaz:
\begin{enumerate}[label=(\alph*)]
\item a mesma equação de pressão~\eqref{eq:pressure-eq} (mesmos pesos,
  mesmo operador $M_q(\alpha)$, mesmas raízes $\alpha=1,\alpha^\ast$);
\item o mesmo índice de cauda que $W$;
\item o mesmo coeficiente de regressão linear $b\approx 1$ contra a
  medida limite $\mu$, a menos de uma dispersão adicional fixa;
\item mas $\Var[\log X \mid \text{primeiros } m \text{ dígitos}] \to
  \Var[\log Y] > 0$ para \emph{todo} $m$, inclusive $m=\infty$: a
  variância condicional nunca colapsa.
\end{enumerate}
\end{proposition}

O item (d) é o cerne: nenhuma estatística calculada a partir de
pressão, índice de cauda, ou coeficiente de regressão contra $\mu$
distingue $W$ de $W\cdot Y$. O nome técnico deste fenômeno é
\emph{endogenia} \cite{AldousBandyopadhyay2005}: o colapso da
variância condicional equivale a $X$ ser mensurável em relação à
$\sigma$-álgebra gerada pelos dígitos (``endógena''); a
Proposição~\ref{prop:gauge} exibe uma solução não-endógena.
Enfatizamos o status lógico desta construção: ela prova que a recursão
\emph{sozinha} não força exatidão; não afirma que um $Y$ não-trivial
de fato ocorre para o objeto aritmético construído a partir de um
único inteiro $v$ (ali, $Y$ teria que ser fabricado a partir do
próprio $v$, não injetado como coordenada auxiliar independente). Estas
duas afirmações nunca devem ser confundidas.

\subsection{Um mecanismo específico descartado}

Um candidato natural para um fator de gauge não-trivial é uma
``memória'' $2$-ádica persistente da raiz sobrevivendo profundamente
na árvore. Isto é falso:

\begin{lemma}[Sem memória $2$-ádica]\label{lem:2adic}
Seja $w$ um filho de $v$ no expoente $a$ na árvore reversa do mapa
clássico ($q=3$), de modo que $3w=2^a v - 1$. Então
\[
  w \equiv -3^{-1} \pmod{2^a},
\]
uma constante universal dependendo só de $a$, independente de $v$.
Consequentemente, um descendente alcançado após um caminho com soma de
expoentes $A$ tem seus $A$ bits mais baixos determinados
\emph{inteiramente pelo caminho}, com informação nula sobre a raiz.
\end{lemma}

\begin{proof}
De $3w=2^av-1$ obtemos $3w\equiv -1\pmod{2^a}$, i.e.\ $w\equiv
-3^{-1}\pmod{2^a}$ já que $3$ é unidade mod $2^a$. Indução ao longo do
caminho, compondo a relação a cada geração, dá a afirmação.
\end{proof}

Verificamos o Lema~\ref{lem:2adic} numericamente para $a=1,\dots,10$
contra mais de $500$ valores de $v$ por expoente, confirmando a
constante universal prevista exatamente em todos os casos (ex.:
$a=5\Rightarrow w\equiv 21\ (\mathrm{mod}\ 32)$ sempre; $a=10\Rightarrow
w\equiv 341\ (\mathrm{mod}\ 1024)$ sempre). A árvore reversa destrói
toda informação $2$-ádica sobre a raiz a cada geração; não existe
coordenada $2$-ádica persistente disponível para correlacionar com o
resíduo $3$-ádico, descartando este mecanismo específico para um fator
de gauge.

\subsection{O que força o colapso, e o que não força}

No modelo probabilístico totalmente idealizado --- ramificação i.i.d.\
com conteúdos de subárvore independentes --- a exatidão \emph{é}
forçada, por dois mecanismos clássicos:

\begin{enumerate}[label=(\alph*)]
\item \textbf{Classificação de pontos fixos de transformações de
  suavização} \cite{DurrettLiggett1983,AlsmeyerBigginsMeiners2012}:
  soluções não-endógenas da equação de ponto fixo de suavização
  associada genericamente têm índice de cauda estritamente menor que a
  endógena (no regime relevante aqui, índice de cauda $1$ com média
  infinita), o que é excluído pelos nossos próprios dados: $\alpha^\ast=2$
  foi medido por três métodos independentes (equação de pressão,
  estimador de Hill, teoria de valores extremos sobre máximos de
  blocos).
\item \textbf{Rigidez de Choquet--Deny} para a fase arquimediana:
  funções harmônicas limitadas de um passeio aleatório num grupo
  abeliano são constantes, o que elimina qualquer coordenada de gauge
  arquimediana sob a hipótese i.i.d.\ completa.
\end{enumerate}

\begin{theorem}[Enunciado preciso da barreira]\label{thm:barrier}
Qualquer argumento que use apenas \textup{(i)} a recursão exata de
pesos da árvore reversa e \textup{(ii)} estatísticas dos dígitos
$3$-ádicos fatora através do modelo de transformação de suavização,
que admite soluções não-endógenas (fatores de gauge harmônicos) com
piso de variância condicional estritamente positivo em toda resolução,
estatisticamente indistinguíveis da solução endógena por pressão,
índice de cauda, ou qualquer marginal que testamos. O ingrediente
faltante que fecha a lacuna \emph{no modelo idealizado} --- e que resta
transferir para o cenário aritmético --- é a quase-independência,
entre subárvores irmãs, dos dígitos $3$-ádicos frescos dos
inteiros-folha, junto com equidistribuição da fase arquimediana ao
longo dos caminhos. O análogo rigoroso conhecido, na direção forward, é
o decaimento de Fourier das variáveis de Syracuse de Tao
(Proposição~\ref{prop:tao117}); a versão árvore-reversa/densidade
permanece aberta e \emph{não} é implicada por ele.
\end{theorem}

Enfatizamos o status epistêmico das duas metades do
Teorema~\ref{thm:barrier}: o resultado de classificação (a) pode ser
enunciado incondicionalmente como ``maquinaria padrão sob [hipóteses],
verificação detalhada pendente'', em vez de um teorema que
re-derivamos independentemente contra nosso setup exato; a exclusão
$2$-ádica (Lema~\ref{lem:2adic}) e a construção de gauge
(Proposição~\ref{prop:gauge}) são totalmente rigorosas e verificadas
independentemente. Nunca afirmamos que a dispersão residual de gauge
$\sigma_Y$ se anula para o objeto aritmético verdadeiro; só que a
recursão sozinha não pode forçar isso.

\subsection{Regime preciso do expoente de cauda}
\label{sec:endogeny-addendum}

Uma busca literária dirigida (\S\ref{sec:context}) levantou a questão
de se $\alpha^\ast=2$ coincide com o \emph{caso crítico} da teoria
clássica de transformações de suavização multivariadas
\cite{KoleskoMentemeier2015}, caracterizado por $m'(\alpha)=0$ numa
raiz de $m(\alpha)=1$ (um caso de fronteira/degenerado no sentido de
Biggins--Kyprianou). Escrevendo $m(\alpha):=q^{\alpha-1}/(2^\alpha-1)$
(a função de pressão da \S\ref{sec:pressure} em $q=3$), calculamos
\[
  m'(1) = \ln 3 - 2\ln 2\cdot\frac{2^1}{2^1-1}
    \quad\text{explicitamente } \approx -0{,}288 < 0, \qquad
  m'(2) = \ln 3 - \tfrac{4}{3}\ln 2 \approx +0{,}174 > 0.
\]
Como $m$ é log-convexa (como soma, após limpar denominadores, de
exponenciais em $\alpha$), duas raízes distintas de $m(\alpha)=1$
necessariamente ladeiam o mínimo único de $m$, logo \emph{nenhuma}
delas pode ser raiz crítica (dupla); criticalidade exigiria que as duas
raízes se fundissem, o que não acontece em $q=3$. A classificação
correta de $\alpha^\ast=2$ é portanto como \emph{segunda raiz com
$m'>0$} --- um regime de cauda regularmente variável do tipo
Goldie/Guivar'ch \cite{Goldie1991,BuraczewskiDamekMikosch2016}, não o
caso crítico. Esta distinção importa operacionalmente: o teorema de
unicidade do caso crítico de \cite{KoleskoMentemeier2015}
\emph{pressupõe} exatamente a independência entre cópias de subárvore
que é a lacuna aritmética do Teorema~\ref{thm:barrier} --- então, mesmo
deixando de lado o descompasso de regime, ele não poderia ter
fornecido uma ferramenta nova para fechar essa lacuna. Registramos,
ligando isso à \S\ref{sec:freezing}, que "segunda raiz com $m'>0$" é
exatamente a condição de congelamento da
Proposição~\ref{prop:always-frozen}: o nome do regime
Goldie/Guivar'ch classifica corretamente a raiz algébrica
$\alpha^\ast=2$ de $m(\alpha)=1$, mas por si só não diz nada sobre se
essa raiz governa a cauda \emph{quenched} da nossa recursão aritmética
específica --- que é exatamente a pergunta que a
Proposição~\ref{prop:always-frozen} e a Conjectura~\ref{conj:tail-index}
endereçam diretamente.

\section{Calibração empírica do acoplamento residual}
\label{sec:calibration}

A Proposição~\ref{prop:gauge} mostra que a lacuna é real em princípio;
não quantifica quão grande é o acoplamento residual entre subárvores
irmãs de fato para a árvore aritmética verdadeira. Desenhamos um
experimento de controle de três braços para calibrar isso diretamente.

\subsection{Desenho}

\begin{enumerate}[label=Braço~\arabic*.]
\item \textbf{Controle sintético i.i.d.}: um modelo de Monte Carlo com
  os pesos de ramo exatos $q\cdot 2^{-a}$ mas conteúdos de subárvore
  genuinamente independentes, truncado por tamanho populacional total
  (não profundidade) para espelhar o procedimento de enumeração
  verdadeiro.
\item \textbf{Controle sintético acoplado}: o mesmo modelo, mas com os
  \emph{conteúdos} de subárvore irmãs acoplados em intensidade
  ajustável $\rho\in[0,1]$ via replicação de família inteira endereçada
  por um hash de (fase, id de réplica, tag), construído de modo que
  todas as marginais de subárvore única permaneçam exatamente
  inalteradas para qualquer $\rho$.
\item \textbf{Árvore aritmética real}: enumeração exata da árvore
  reversa verdadeira, remedida na mesma profundidade de truncamento.
\end{enumerate}

O excesso de variância residual, $(\text{Braço 3}) - (\text{Braço 1})$,
comparado contra a curva de calibração $\text{piso}(\rho)$ construída
a partir do Braço 2, converte uma afirmação qualitativa (``a lacuna
existe'') numa cota quantitativa sobre o acoplamento efetivo
$\rho_{\mathrm{eff}}$ que seria necessário no modelo sintético para
reproduzir o resíduo observado na árvore real.

\subsection{Resultado}

Ao longo de toda a grade de truncamento $m=8,\dots,29$, os Braços 1 e 3
mostraram variâncias residuais quase idênticas, sem sinal de
acoplamento positivo. Uma medição final de inclinação com bootstrap no
truncamento mais profundo e melhor calibrado ($m=20$, $4000$ raízes
replicadas, $\rho\in\{0,0{,}05,0{,}1,0{,}2\}$) dá:
\begin{equation}\label{eq:rho-eff}
  \rho_{\mathrm{eff}} \lesssim 0{,}06 \quad (\text{IC } 95\%, \ m=20).
\end{equation}
O excesso (real $-$ sintético) converge monotonicamente a zero com o
aumento da profundidade de truncamento, o sinal oposto ao que um
acoplamento aritmético positivo genuíno produziria, reforçando em vez
de apenas deixar de contradizer a conclusão. Nenhum acoplamento
aritmético positivo entre subárvores irmãs foi detectado; o
acoplamento efetivo é limitado a aproximadamente $6\%$ do acoplamento
máximo explorado no controle sintético. A equação~\eqref{eq:rho-eff}
converte a lacuna puramente teórica do Teorema~\ref{thm:barrier} numa
cota empírica quantitativa e falsificável do seu tamanho.

\section{Um resultado empírico incondicional: resolvendo a disputa
Kontorovich--Lagarias vs.\ Volkov para $5x+1$}
\label{sec:kl}

Independentemente do programa teórico acima, relatamos um resultado
empírico totalmente incondicional de interesse direto.

\subsection{A disputa}

Para a árvore reversa de $T_5$, Kontorovich--Lagarias
\cite{KontorovichLagarias2009} preveem um expoente de contagem
$\eta_{5,BP}\approx 0{,}650919$ (idêntico, como notado na
Observação~\ref{rem:novelty-109}, à nossa raiz da equação de pressão
$\alpha_-(5)$), enquanto um modelo estocástico concorrente devido a
Volkov, discutido na mesma referência, prevê
$\eta^\ast_{5,BP}\approx 0{,}678$. Os autores de
\cite{KontorovichLagarias2009} eles mesmos sinalizam isso como uma
disputa em aberto, não resolvida por falta de dados suficientemente
poderosos; as duas previsões diferem por apenas $\Delta=0{,}027$.

\subsection{Método}

Implementamos a árvore reversa exata de $T_5$ do zero (a regra de
admissibilidade da \S\ref{sec:setup} especializada a $q=5$, $d=4$: um
filho no expoente $a$ existe a partir do pai $u$ sse $a\equiv
A_0(u\bmod 5)\pmod 4$, com $A_0=\{1\!\mapsto\!4,\allowbreak
2\!\mapsto\!3,\allowbreak 3\!\mapsto\!1,\allowbreak 4\!\mapsto\!2\}$;
nós com $u\equiv 0\pmod 5$ são estéreis). Diferente de $q=3$, não há
condição de paridade além da integralidade mod $5$.

Um primeiro piloto ($60$ raízes) já deu uma inclinação de $0{,}6433$
(IC 95\% $[0{,}6327,0{,}6531]$), excluindo $0{,}678$ por cerca de
$6{,}4$ erros-padrão. Uma execução de produção maior ($n=300$ raízes)
inicialmente amostrou raízes próximas demais do checkpoint inferior da
janela de medição, contaminando o buffer de truncamento; corrigindo
isso (raízes restritas bem abaixo da janela) deu uma inclinação de
$0{,}63801$ (IC 95\% $[0{,}63159, 0{,}64410]$) --- que, neste tamanho
de amostra, excluiu \emph{ambos} os valores candidatos ($0{,}678$ por
$12{,}6$ erros-padrão, mas também $0{,}650919$ por $4{,}1$
erros-padrão), um modo de falha antecipado previamente: uma vez que o
erro-padrão cai abaixo do viés de truncamento residual, um intervalo de
confiança ingênuo sobre o estimador enviesado exclui tudo, o que não é
evidência contra a previsão correta.

Por isso reescrevemos a enumeração para rastrear, numa única passada,
o máximo-de-caminho corrente em cada nó, dando contagens simultâneas
em buffers de truncamento $10^9$ a $10^{13}$ (validado byte a byte
contra o método anterior buffer-a-buffer em quatro raízes de teste). A
curva média agregada entre raízes decai geometricamente:
\[
  0{,}60049\ (10^9) \to 0{,}62387\ (10^{10}) \to 0{,}63261\ (10^{11}) \to
  0{,}63650\ (10^{12}) \to 0{,}63801\ (10^{13}),
\]
com incrementos $0{,}0234, 0{,}0087, 0{,}0039, 0{,}0015$ encolhendo por
um fator de aproximadamente $0{,}4$--$0{,}5$ por década de truncamento
--- assinatura de um viés de truncamento genuíno e extrapolável, não
ruído.

\subsection{Resultado}

A extrapolação de Aitken $\Delta^2$ desta sequência para truncamento
infinito dá
\begin{equation}\label{eq:kl-result}
  \hat\eta_5 = 0{,}639, \qquad \text{IC } 95\%\ \text{(bootstrap sobre
  raízes)} = [0{,}633,\, 0{,}645].
\end{equation}
Isso exclui a previsão de Volkov ($0{,}678$) com margem ampla (mais de
dez erros-padrão). O pequeno resíduo até o valor de
Kontorovich--Lagarias ($0{,}650919$, cerca de $0{,}012$ acima do nosso
limite superior de confiança) tem uma assinatura específica e
diagnosticável em vez de ser uma discrepância inexplicada: o painel de
inclinação-por-década \emph{dentro} da janela de medição fixa
($10^4\!\to\!10^5$ até $10^7\!\to\!10^8$) ainda está subindo
monotonicamente na última década testada
($0{,}6021\to0{,}6296\to0{,}6432\to0{,}6460$, sem platô), aproximando-se
de $0{,}6509$ monotonicamente sem tê-lo alcançado --- a assinatura de
\emph{pré-assintótica de janela fixa} (a janela ainda não é
``profunda o suficiente''), não de viés de truncamento não corrigido
(já tratado pela extrapolação de Richardson acima).

\begin{theorem}[Resolução empírica da disputa Kontorovich--Lagarias
vs.\ Volkov]\label{thm:kl}
O expoente de contagem da árvore reversa de $5x+1$, medido por
enumeração exata com correção de truncamento por extrapolação de
Richardson, é $0{,}639$ (IC 95\% $[0{,}633,0{,}645]$). Isso exclui a
previsão concorrente de Volkov ($0{,}678$) com alta confiança e é
consistente com a previsão de Kontorovich--Lagarias ($0{,}650919$) a
menos de um efeito de pré-assintótica de janela finita, medido e
explicado.
\end{theorem}

Este resultado se sustenta inteiramente por si só, independente do
programa teórico das \S\S\ref{sec:endogeny}--\ref{sec:lemmas};
destacamo-lo cedo no argumento geral do paper porque é a contribuição
menos dependente de enquadramento que relatamos.

\section{Teste computacional direto da Weak Covering Conjecture}
\label{sec:wcc}

\subsection{A conjectura de Wirsching}

Wirsching \cite{Wirsching1998} identifica a seguinte propriedade
combinatória de cobertura como suficiente para densidade positiva
uniforme de predecessores (seu Corolário~II.5.8 e Capítulo V). Defina
\[
  R_{j,k} := \Bigl\{ \textstyle\sum_{i=0}^j 2^{\alpha_i} 3^i \;:\;
  j+k\ge \alpha_0>\alpha_1>\cdots>\alpha_j\ge 0 \Bigr\} \subset
  \mathbb{Z}_3^\ast, \qquad |R_{j,k}| = \binom{j+k+1}{j+1}.
\]
Todo elemento de $R_{j,k}$ é automaticamente unidade mod $3$ (seu
termo de maior ordem não é divisível por $3$, todos os demais são),
então ``cobrir $\mathbb{Z}_3^\ast$ mod $3^\ell$'' significa atingir
todas as $2\cdot 3^{\ell-1}$ classes de resíduo invertíveis.

\begin{conjecture}[Weak Covering Conjecture; Wirsching 1998,
Conj.~3.9]\label{conj:wcc}
Existe uma função sub-exponencial $K(\ell)>0$ (i.e.\
$\lim_\ell K(\ell)e^{-\gamma\ell}=0$ para todo $\gamma>0$) tal que,
para todo $j,\ell$: $|R_{j-1,j}|\ge K(\ell)\cdot 3^\ell \Rightarrow
R_{j-1,j}$ cobre $\mathbb{Z}_3^\ast \bmod 3^\ell$.
\end{conjecture}

Para $j\ge \ell$, termos de índice $i\ge\ell$ se anulam mod $3^\ell$,
dando a redução útil $R_{j-1,j}\bmod 3^\ell = 2^{j-\ell}\cdot
R_{\ell-1,j}\bmod 3^\ell$, que reduz o custo de enumeração de
$\binom{2j}{j}$ para $\binom{j+\ell}{\ell}$.

\subsection{Implementação e validação}

Implementamos uma programação dinâmica de bitset com rotação cíclica:
processando expoentes candidatos em ordem decrescente, mantemos, para
cada contagem $c$ de termos já escolhidos, o conjunto de somas parciais
alcançáveis mod $3^\ell$ como um único inteiro bitset em Python;
escolher o expoente $v$ na posição $c$ é uma rotação cíclica por
$2^v\cdot 3^c\bmod 3^\ell$, e a contagem de bits do bitset final dá o
tamanho da imagem diretamente. A implementação bateu exatamente com uma
tabela de referência calculada independentemente ($j^\ast(\ell)$ para
$\ell=1,\dots,7$: $1,4,6,7,9,10,11$) e uma checagem pontual
($\ell=2,j=2$: imagem $=\{1,2,5,7\}\bmod 9$, faltando $\{4,8\}$).

\subsection{Resultados}

Calculamos $j^\ast(\ell)$, o menor $j$ para o qual $R_{j-1,j}$ cobre,
para $\ell=1,\dots,20$ (o custo cresceu por um fator $\approx 3{,}3$
por passo perto do fim; $\ell=20$ sozinho levou $\approx 27$ minutos;
extensão adicional em Python puro foi julgada não compensar).
Escrevendo $e(\ell):=j^\ast(\ell)-\log_4 3\cdot\ell$ para o excesso
sobre o limiar casado por entropia $\log_4 3\approx 0{,}7925$:

\begin{center}
\begin{tabular}{c|c|c|c}
$\ell$ & $j^\ast$ & $j^\ast/\ell$ & $e(\ell)$ \\ \hline
1 & 1 & 1{,}0000 & 0{,}208 \\
5 & 9 & 1{,}8000 & 5{,}038 \\
10 & 15 & 1{,}5000 & 7{,}075 \\
15 & 20 & 1{,}3333 & 8{,}113 \\
16 & 20 & 1{,}2500 & 7{,}320 \\
17 & 21 & 1{,}2353 & 7{,}528 \\
18 & 22 & 1{,}2222 & 7{,}735 \\
19 & 23 & 1{,}2105 & 7{,}943 \\
20 & 24 & 1{,}2000 & 8{,}150
\end{tabular}
\end{center}

(A tabela completa com os $20$ valores está no registro computacional
que acompanha o paper.) $j^\ast(\ell)$ existe e é finito em toda a
faixa testada. Uma leitura anterior dos primeiros $17$ pontos sugeria
uma inclinação local desacelerando monotonicamente na direção do
limiar $\log_4 3$, o que os três novos pontos em $\ell=18,19,20$
(todos os incrementos exatamente $+1$, sem novos platôs) revelaram ser
artefato de um platô isolado em $\ell=16$ atravessando a janela de
medição deslizante: sem ele, a inclinação local oscila em
$0{,}74$--$0{,}86$, cavalgando $\log_4 3$ de ambos os lados em vez de
convergir de forma limpa.

Comparação de modelos na cauda $\ell=10$--$20$ (crescimento constante
vs.\ logarítmico vs.\ raiz-quadrada vs.\ linear-lento de $e(\ell)$, por
AIC) agora desfavorece a estabilização pura
($\Delta\mathrm{AIC}\approx 5$), consistente com um teste de
frequência de platôs (apenas $1$ platô observado em $19$ incrementos
contra $\approx 3{,}9$ esperados sob estabilização na taxa $\log_4 3$;
binomial $p\approx 0{,}072$, marginal mas na mesma direção). Os quatro
modelos de crescimento lento permanecem estatisticamente
indistinguíveis entre si nesta faixa ($\Delta\mathrm{AIC}<0{,}5$); eles
só se separariam por $\ge 1$ unidade de $j^\ast$ perto de $\ell\approx
40$, fora do alcance desta abordagem algorítmica.

\begin{theorem}[Teste computacional direto da Weak Covering
Conjecture]\label{thm:wcc}
$j^\ast(\ell)$ existe e é finito para todo $\ell\le 20$ (o primeiro
teste computacional direto, até onde sabemos, do objeto exato da
Conjectura~\ref{conj:wcc}). O excesso $e(\ell)$ favorece
qualitativamente crescimento lento ilimitado sobre estabilização
(duas estatísticas independentes, $\Delta\mathrm{AIC}$ e frequência de
platôs, na mesma direção mas não individualmente decisivas),
consistente com a forma fraca (Conjectura~\ref{conj:wcc}) e
desfavorecendo a forma forte correspondente com $K(\ell)$ limitado.
\end{theorem}

\section{Três regimes de precisão para a extensão bivariada de Tao}
\label{sec:regimes}

\subsection{Por que o argumento bivariado ingênuo falha}

Uma estratégia natural para fechar a lacuna do Teorema~\ref{thm:barrier}
é aplicar a Proposição~\ref{prop:tao117} duas vezes, uma por folha
irmã, e combinar via alguma forma de independência condicional dado o
ancestral comum. Uma leitura cuidadosa do mecanismo da prova da
Seção~7 de Tao (caracteres $\chi_\xi$, pares de valuações i.i.d.\
distribuídas como Pascal, uma fase evoluindo multiplicativamente por
$\log 9/\log 2$, um ``conjunto preto'' de triângulos bem separados, um
processo de renovação bidimensional com monotonicidade) mostra que
esse mecanismo opera inteiramente sobre
$\mathrm{Syrac}(\mathbb{Z}/3^n\mathbb{Z})$ construída a partir de
dígitos geométricos \emph{independentes} por definição
\eqref{eq:syracuse-def}; a ligação com inteiros reais é feita em outro
lugar de \cite{Tao2022} (via uma aproximação em variação total válida
só para profundidade $n\lesssim\log N$).

Para o nosso problema --- um único inteiro fixo $v$, duas folhas irmãs
$w_1(v),w_2(v)$ --- as duas folhas \emph{não} são duas amostras
i.i.d.; são duas funções determinísticas diferentes da \emph{mesma}
variável $v$.

\begin{proposition}[Dependência perfeita em profundidade
fixa]\label{prop:fixed-pair}
Para um par fixo de caminhos irmãos de profundidade $D$, a congruência
de existência pina $v$ numa única classe mod $3^D$: escrevendo
$v=v_0+3^D t$, as duas folhas viram funções afins da mesma variável
livre $t$, com multiplicadores unitários: $w_i=w_i(v_0)+2^{A_i}t$.
Consequentemente o par $(w_1,w_2)\bmod 3^\ell$ vive numa reta de
tamanho $3^\ell$ dentro de $(\mathbb{Z}/3^\ell\mathbb{Z})^2$, para todo
$\ell\ge 1$ --- dependência determinística perfeita em toda precisão,
com uma reta inteira de frequências ressonantes onde a correlação tem
módulo exatamente $1$.
\end{proposition}

Consequentemente, ``aplicar a Proposição~\ref{prop:tao117} duas vezes,
assumindo independência condicional dado o ancestral'' não é apenas
circular: é falso como enunciado, em qualquer precisão.

\subsection{A reformulação correta, não circular}

A quantidade que de fato importa --- decorrelação de \emph{agregados
de subárvore}, não folha-a-folha --- é um enunciado de segundo
momento, que se expande numa soma sobre \emph{pares de caminhos}:
\[
  \mathbb{E}_v[Z_1 Z_2] \;\propto\; \#\{(a,a') : a_1\in B_1,\ a'_1\in
  B_2,\ \mathrm{Syrac}(a)\equiv\mathrm{Syrac}(a') \bmod 3^D\},
\]
que por inversão de Fourier dá $\Cov \propto \sum_{\xi\ne 0}
S_1(\xi)S_2(\xi)^\ast$, com $S_i$ somas de caracteres condicionadas no
primeiro passo. Aqui não há circularidade: os dois índices de caminho
são independentes pela tautologia de expandir um quadrado, não por
uma hipótese aritmética.

\subsection{Os três regimes}

\begin{enumerate}
\item \textbf{Par fixo, qualquer $\ell$}: dependência perfeita
  (Proposição~\ref{prop:fixed-pair}); nenhum decaimento é possível, e
  decorrelação só pode vir de média sobre caminhos, nunca
  folha-a-folha.
\item \textbf{$\ell = O(\log D)$}: \emph{demonstrável}, com a
  maquinaria da Seção~7 de Tao essencialmente inalterada. O orçamento
  $|\rho(\xi)|\le C_A D^{-A}$ fecha a soma sobre $3^\ell-1$ frequências
  desde que $3^\ell \ll D^{2A}$, i.e.\ módulos polinomiais em $D$ ---
  um lema honesto e alcançável, embora sua formulação mais natural
  acabe sendo falsa como originalmente esboçada (\S\ref{sec:lemmas}).
\item \textbf{$\ell \asymp c\cdot D$} (onde vivem a barreira de
  endogenia, os dígitos frescos, e a Weak Covering Conjecture de
  Wirsching): isso exige decaimento de Fourier \emph{exponencial}
  (power-saving), uniforme em $\xi$; a Proposição~\ref{prop:tao117}
  só dá decaimento super-polinomial, e o método da Seção~7 de Tao
  estruturalmente não pode dar mais, porque as perdas no argumento do
  conjunto preto estão amarradas a propriedades diofantinas de
  $\log 3/\log 2$ via a inclinação $\log 9/\log 2$.
\end{enumerate}

\begin{theorem}[O regime 3 é um problema aberto reconhecido]
\label{thm:regime3}
O regime 3 acima é equivalente, em dificuldade, a (i) correlações de
funções digitais nas bases $2$ e $3$ \cite{Spiegelhofer2021}, (ii)
rigidez efetiva para estruturas $\times 2,\times 3$-invariantes no
sentido de Furstenberg (o par $(w_1,w_2)$ é $v$ observado através de
dois elementos do semigrupo afim gerado por $x\mapsto(2^ax-1)/3$),
aberto na generalidade necessária aqui, e (iii) decaimento de Fourier
exponencial para medidas autossimilares $3$-ádicas do tipo
Breuillard--Varjú, conhecido só em casos algébricos especiais. Nenhum
destes fornece uma ferramenta utilizável hoje.
\end{theorem}

\section{Triangulação: três formulações independentes da mesma
obstrução}
\label{sec:triangulation}

\subsection{$\beta=1$ (Tao 2020) é a Weak Covering Conjecture
(Wirsching 1998)}

Tao \cite{Tao2020} define $c_n := \inf_{b\in\mathbb{Z}/3^n\mathbb{Z},\,
3\nmid b} \Pr(\mathrm{Syrac}(\mathbb{Z}/3^n\mathbb{Z})=b)$ e propõe a
conjectura $\beta=1$: $c_n = 3^{-n+o(n)}$ (``nenhuma classe de resíduo
tem probabilidade muito menor que a média''), provando
condicionalmente que $\beta=1$ implica densidade de Collatz $\gg
x^{1-o(1)}$ (quase total), e notando que a melhor cota incondicional
conhecida é o $\gamma=0{,}84$ de Krasikov--Lagarias
\cite{KrasikovLagarias2003}.

\begin{proposition}[Identidade estrutural]\label{prop:beta-wcc}
Após um twist de unidade, a variável aleatória de Syracuse de Tao e os
conjuntos $R_{j,k}$ de Wirsching são o mesmo objeto algébrico (somas
de potências de $2$ e $3$ com expoentes estritamente decrescentes).
Consequentemente, a Conjectura~\ref{conj:wcc} ($K(\ell)$
sub-exponencial) implica $\beta=1$ (uma conta de entropia exata:
$j=\log_4 3\cdot \ell + o(\ell)$ dá $c_\ell \gtrsim 3^{-\ell(1+o(1))}$);
reciprocamente $\beta=1$ sozinho dá apenas uma forma exponencialmente
enfraquecida da Conjectura~\ref{conj:wcc}. As duas são formulações
irmãs na mesma escala exata de entropia --- o limiar
$\log_4 3\approx 0{,}7925$ da \S\ref{sec:wcc} é a mesma identidade
$4^j$ vs.\ $3^\ell$ nos dois lados, não uma coincidência.
\end{proposition}

\begin{corollary}\label{cor:e-ell-beta}
O excesso $e(\ell)=j^\ast(\ell)-\log_4 3\cdot \ell$ medido na
\S\ref{sec:wcc} é literalmente o termo $o(n)$ de $\beta=1$
($c_\ell\gtrsim 3^{-\ell}\cdot 4^{-e(\ell)}$): nossa leitura
qualitativa de $e(\ell)$ (crescimento sub-linear favorecido) é
evidência computacional direta a favor de $\beta=1$, não apenas da
Conjectura~\ref{conj:wcc}.
\end{corollary}

A abordagem de Tao de $2020$ também mostra que a rota de
endogenia/decorrelação da \S\ref{sec:endogeny} não é o único caminho:
ela substitui ``quase-independência entre subárvores'' por
``separação determinística de pré-imagens (fácil) mais controle
marginal $L^\infty$ ($\beta=1$, onde mora toda a dificuldade)''. Isto
confirma que existe uma rota que evita a decorrelação por completo, e
que ela bate na \emph{mesma} parede em forma marginal.

\subsection{A equação funcional de Wirsching de 2003 e a função de
Fabius em base $3$}

Uma rota diferente e mais refinada em direção ao mesmo objetivo é
Wirsching~\cite{Wirsching2003}, cujo objeto combinatório central, as
funções ``Elka'' $e_\ell(k,a):=|E_{\ell,k}(a)|$, contam caminhos de
comprimento $k+\ell$ no grafo de Collatz terminando em $a$ --- o mesmo
objeto de contagem que sustenta todo o nosso programa
$G(v)/\mathrm{Syrac}$, agora visto no espaço de estados
$\mathbb{X}=\mathbb{I}_3\times \mathbb{Z}_3^\times$: o fator
$\mathbb{Z}_3^\times$ carrega a variável $3$-ádica (habitat da medida
de Syracuse), enquanto o fator $\mathbb{I}_3\cong\{0,1,2\}^{\mathbb{N}}$
carrega a contagem normalizada de passos $k/3^\ell$, convergindo para
uma densidade $\varphi$.

\begin{proposition}[Função de Fabius em base $3$]\label{prop:fabius}
$\varphi$ é a densidade de $X = \sum_{j\ge 1} 2U_j\cdot 3^{-j}$, $U_j$
i.i.d.\ uniformes em $[0,1]$ --- o análogo em base $3$ da função de
Fabius (que satisfaz $f'(x)=2f(2x)$; aqui $\varphi'(x) =
\tfrac{9}{2}\varphi(3x)$ em $[0,2/3]$), membro da família de
``funções atômicas'' de Rvachev. É o único ponto fixo em
$L^1([0,1])$ do operador de médias $W_3 f(x) := \tfrac{3}{2}
\int_{3x-2}^{3x} f(t)\,dt$, $C^\infty$, polinomial por partes fora do
conjunto de Cantor padrão, com convergência geométrica certificada
$\|f_n-\varphi\|_1\le 2^{-n+1}\|f_1-f_0\|_1$.
\end{proposition}

A cadeia de implicações de Wirsching reduz densidade positiva uniforme
de predecessores a uma sequência de três conjecturas, a mais concreta
das quais (Conjectura~3, uma desigualdade sobre
$\liminf \varphi(z_\ell)/\varphi_0(z_\ell)$ ao longo de sequências na
janela do teorema central do limite $|\ell-k_\ell|\le \delta\sqrt\ell$)
testamos numericamente.

\subsection{Teste numérico certificado da Conjectura 3 de Wirsching}

Usando a autossimilaridade exata $X\overset{d}{=}(2U+X)/3$, os momentos
$M_i=\mathbb{E}[X^i]$ satisfazem a recursão racional exata
\[
  M_i = \frac{1}{3^i-1}\sum_{k=1}^i \binom{i}{k}\frac{2^k}{k+1} M_{i-k},
\]
e, combinado com uma redução por antiderivadas expressando $\varphi$
em argumentos de cauda extrema em termos de somas binomiais exatas de
momentos (em vez de iterar $W_3$, cujo certificado de erro em $L^1$ é
fraco demais para controlar valores pontuais na escala relevante),
obtemos uma avaliação de $\varphi$ com erro de truncamento zero em
profundidade $\ell$ até $\ell=500$ (trabalhando o tempo todo em
espaço-logarítmico com precisão de $60$ dígitos; erro certificado
$\le 10^{-8}$).

\begin{theorem}[Teste numérico da Conjectura 3 de Wirsching]
\label{thm:conjecture3}
A razão decisiva $L_\ell := 3^{1-\ell}\varphi(3x_\ell^+)/\varphi(x_\ell^+)$,
medida para $\ell=250,\dots,500$ na janela do TCL ($u\in[-2,2]$),
converge ao valor previsto $2/3$ com déficit
$\ell\cdot(2/3-L_\ell) \to 0{,}580\pm 0{,}001$ --- um coeficiente
reproduzido independentemente pela própria assintótica $\varphi_0$ de
Berg--Krüppel. A quantidade $\ln(\varphi/\varphi_0)$ é crescente,
côncava, e uniforme em $u\in[-2,2]$ (dispersão $<10^{-4}$),
consistente com um limite finito
$L=-0{,}619\pm0{,}001\,(\text{estat.})$, com incerteza sistemática de
forma de cauda de $\pm 0{,}015$, dando $c=e^L\approx 0{,}538$ (faixa
honesta $0{,}53$--$0{,}55$), com melhor ajuste por um termo de cauda
$0{,}79/\sqrt\ell$; nenhuma modulação log-periódica é detectada. Isto
suporta a Conjectura~3 (a mais concreta da cadeia de Wirsching) com
erro numérico certificado, na faixa testada; as Conjecturas~1 e~2
acima dela permanecem abertas.
\end{theorem}

Juntos, o Teorema~\ref{thm:wcc}, a Proposição~\ref{prop:beta-wcc}, e o
Teorema~\ref{thm:conjecture3} dão três ``roupagens'' formais
independentes --- contagem de inteiros (Wirsching 1998), caracteres de
Fourier (Tao 2020), uma equação funcional real (Wirsching 2003) --- da
mesma janela crítica, nenhuma resolvida.

\section{Dois lemas candidatos, executados até o fim}
\label{sec:lemmas}

As \S\S\ref{sec:regimes}--\ref{sec:triangulation} identificam dois
próximos passos naturais: um ``lema do regime 2'' (decorrelação de
agregados de subárvore em precisão logarítmica) e um ``lema de redução
Z-number'' (formalizando a analogia com os métodos de
Littlewood--Offord/conjuntos de Bohr de \cite{Tao2011}). Executamos os
dois até o fim; nenhum sobrevive como originalmente esboçado, mas
ambos rendem resultados genuínos e citáveis.

\subsection{O lema do regime 2: refutação, teorema de estrutura, e uma
refutação computacional de sua condição suficiente natural}

O orçamento originalmente esboçado, $|\rho(\xi)|\le C_A D^{-A}$
\emph{uniformemente em $\xi$}, contradiz a Proposição~\ref{prop:tao123}
(uma identidade exata, não uma estimativa com perda): para qualquer
frequência $\xi$ de nível primitivo $\ell'$ (i.e.\
$\xi=3^{\ell-\ell'}\xi_0$, $3\nmid\xi_0$) e qualquer profundidade
$D\ge\ell'$, $\mathbb{E}[e(\xi F_D/3^\ell)] = \varphi_{\ell'}(\xi_0)$ é
\emph{independente de $D$}: profundidade além de $\ell'$ é um
parâmetro mudo, porque a ultrametricidade $3$-ádica trunca a soma
relevante às $\ell$ coordenadas mais grosseiras
(Proposição~\ref{prop:tao123}). Um exemplo fechado em $\ell'=1$:
$\mathrm{Syrac}_1 = 2^{-a}\bmod 3 \in\{1,2\}$ com probabilidades
$1/3,2/3$; $|\varphi(1)|=1/\sqrt3\approx 0{,}577$ para \emph{todo} $D$,
sem decaimento algum. O orçamento correto por frequência é
$C_A\ell'^{-A}$ (Proposição~\ref{prop:tao117}, que exige
$3\nmid\xi$), não $C_AD^{-A}$; mesmo com essa correção, o orçamento
total $\ell^\infty$ somado sobre frequências, $\sum_{\ell'\le \ell}
3^{\ell'}(C_A\ell'^{-A})^2 \asymp 3^\ell \ell^{-2A}$, nunca tende a
$0$: fechar por contagem de frequências exigiria cancelamento de raiz
quadrada, decaimento $\ell^{-\ell'/2-\varepsilon}$ --- o Regime~3
ressurgindo dentro do que deveria ser o regime fácil.

O que sobrevive é um teorema de estrutura genuíno. Com $w_2 =
2^\Delta w_1 + c_\Delta$, $c_\Delta=(2^\Delta-1)/3$ (um multiplicador
unitário mod $3^\ell$) relacionando folhas irmãs admissíveis com lacuna
de expoente $\Delta$:

\begin{lemma}[Identidade espectral da covariância]\label{lem:cov-spectral}
Para observáveis de agregado de subárvore $Z_i:=N_{\eta_i}^{(m)}(w_i(v))$
(populações ponderadas truncadas em precisão $m$), com $v$ uniforme
nas classes admissíveis mod $3^{m+1}$,
\[
  \Cov(Z_1,Z_2) = 3^{-2m}\sum_{\xi\ne 0 \bmod 3^m}
  e(-\xi c_\Delta/3^m)\, S_1(-2^\Delta\xi)\, S_2(\xi),
\]
por ortogonalidade de caracteres e a relação afim entre irmãs. Os dois
índices de soma são independentes pela tautologia de expandir o
quadrado (não uma hipótese aritmética), mas a dependência através do
$v$ comum não é assim cancelada: ela sobrevive inteiramente como o
emparelhamento diagonal de frequências $\xi_1=-2^\Delta\xi_2$.
\end{lemma}

\begin{proposition}[Endogenia grosseira exata]\label{prop:exact-endogeny}
Para populações completas ($\eta_i\equiv 1$) e precisão $\ell=1$:
$\Corr(Z_1,Z_2)=+1$ se $\Delta\equiv 0\pmod 6$, e $=-1/2$ caso
contrário --- independente da profundidade $D$ e de qualquer
truncamento.
\end{proposition}

A Proposição~\ref{prop:exact-endogeny} converte a barreira de
endogenia do Teorema~\ref{thm:barrier} de um argumento qualitativo em
forma fechada: agregar folhas \emph{não} as decorrelaciona; a
agregação projeta na $\sigma$-álgebra grosseira gerada por
$w_i\bmod 3^\ell$, onde a relação afim entre irmãos do regime de par
fixo (\S\ref{sec:regimes}, Regime~1) reaparece intacta em vez de
desaparecer. Este é um resultado positivo com o sinal oposto ao que se
esperava: o que é demonstrável é a persistência da correlação
grosseira, não seu desaparecimento.

Decompondo a soma do Lema~\ref{lem:cov-spectral} por nível primitivo e
aplicando Cauchy--Schwarz mais Parseval dá $\Cov =
C_{\mathrm{exato}}(\ell_0;\Delta) + \mathrm{Err}$, com
$C_{\mathrm{exato}}$ exato, finito, e independente de $D$
(Proposição~\ref{prop:exact-endogeny}), e
\[
  \frac{|\mathrm{Err}|}{\mathbb{E}[Z_1]\mathbb{E}[Z_2]} \;\le\;
  K_m - K_{\ell_0}, \qquad K_\ell := 3^\ell \sum_y \mu_\ell(y)^2,
\]
a \emph{massa de colisão} normalizada da medida de Syracuse $\mu$. Um
lema condicional seguiria se $K_\infty := \lim_\ell K_\ell < \infty$
(equivalentemente, a densidade $f=d\mu/d(\text{Haar})$ está em
$L^2(\mathbb{Z}_3)$): a covariância fina, descontada da componente
grosseira exata, seria então $\le K_\infty - K_{\ell_0}\to 0$,
uniformemente em $m,D,\Delta$.

\begin{theorem}[Refutação computacional da hipótese $L^2$]
\label{thm:l2-refutation}
$K_\ell$ foi calculado exatamente via a recursão
\begin{equation}\label{eq:mu-recursion}
  \mu_n(y) = \tfrac12\,\nu(2y\bmod 3^n) + \tfrac12\,\mu_n(2y\bmod 3^n),
\end{equation}
onde $\nu$ é a lei de $1+3F_{n-1}$ (derivada abaixo), resolvida como
série geométrica truncada ao longo da órbita cíclica de multiplicação
por $2$ mod $3^n$ (uma única órbita, já que $2$ é raiz primitiva mod
$3^n$ para todo $n$). Para $\ell=1,\dots,17$, $K_\ell$ cresce de forma
limpa e consistentemente \emph{linear}, com incrementos convergindo
monotonicamente a $\approx 0{,}47$, sem sinal de saturação ($K_1=5/3$
exatamente, batendo com um cálculo manual a partir de
$\Pr(\mathrm{Syrac}_1=1)=1/3$, $\Pr(\mathrm{Syrac}_1=2)=2/3$). Logo
$K_\infty=\infty$: a hipótese $L^2$ falha.
\end{theorem}

\begin{proof}[Derivação da recursão~\eqref{eq:mu-recursion}]
Pela Proposição~\ref{prop:tao123}, $F_n(a) =
2^{-a_1}\bigl(1+3F_{n-1}(a_2,\dots,a_n)\bigr)\bmod 3^n$. Escreva
$X:=1+3F_{n-1}(a_2,\dots,a_n)$, com lei $\nu$, independente de $a_1$.
Condicionando em $a_1=1$ (probabilidade $1/2$) obtemos $F_n = 2^{-1}X$;
condicionando em $a_1\ge 2$ (probabilidade $1/2$), a ausência de
memória da distribuição geométrica dá $a_1-1\sim\mathrm{Geom}(2)$
novamente, então $F_n = 2^{-1}\bigl(2^{-(a_1-1)}X\bigr)$, e a
quantidade entre parênteses tem exatamente a lei de $F_n$ (uma cópia
independente). Isso dá o ponto fixo autorreferente $\mu_n(y) =
\tfrac12\nu(2y) + \tfrac12\mu_n(2y)$, toda a aritmética mod $3^n$.
\end{proof}

Validamos a recursão --- não apenas seu caso base $\ell=1$, que não
exercitaria a maquinaria recursiva --- contra amostragem direta de
Monte Carlo de $\mathrm{Syrac}$ a partir de sua definição primária
\eqref{eq:syracuse-def} em $\ell=3,4$ ($3\times10^6$ amostras cada),
bin a bin: discrepância máxima $0{,}000311$ ($\ell=3$) e $0{,}000213$
($\ell=4$), ambas dentro do ruído esperado de Monte Carlo
($\approx 1/\sqrt{N}\approx 0{,}000577$).

\begin{remark}[Conexão com o expoente de cauda]
$H$-104 (uma etapa anterior deste projeto) mediu um expoente de cauda
$\alpha\approx 2$ para razões de população de subárvore --- exatamente
o expoente crítico da condição $L^2$ ($\alpha>2\Rightarrow
K_\infty<\infty$; $\alpha<2\Rightarrow K_\ell$ diverge
polinomialmente; $\alpha=2$ é o caso limítrofe). O crescimento linear
limpo do Teorema~\ref{thm:l2-refutation} é consistente com $\alpha$
ligeiramente abaixo de $2$, ou exatamente no limiar com uma correção
logarítmica.
\end{remark}

\begin{corollary}[Uma escada de ``achatamento'' para a medida de
Syracuse]\label{cor:flatness-ladder}
$L^1$ (provado, Prop.~1.14 de Tao: $\mathrm{Osc}_{m,n}\ll_A m^{-A}$,
dando aproximabilidade em variação total por medidas absolutamente
contínuas) $\;\succ\;$ $L^2$ (refutado aqui,
Teorema~\ref{thm:l2-refutation}) $\;\succ\;$ $L^\infty/\beta=1$
(\S\ref{sec:triangulation}, em aberto) --- nem comparável nem
implicado por $L^2$, mas vivendo no mesmo andar ``além da variação
total'' da escada.
\end{corollary}

O lema do regime 2 não é, portanto, um degrau mais fácil rumo ao
Regime~3: é um irmão dele, compartilhando o mesmo ingrediente
faltante, agora com uma refutação computacional direta da condição
suficiente natural em vez de mero silêncio.

\subsection{A redução Z-number: uma dicotomia condicional e um
resultado negativo exato}

O segundo lema candidato formaliza ``falha da Weak Covering Conjecture
em escala $\ell$ com margem $c$ implica uma configuração de
concentração tipo Bohr'', via a técnica de
Littlewood--Offord/produtos de Riesz do post de blog de Tao de 2011
\cite{Tao2011} (não é um paper; sinalizamos essa correção
bibliográfica explicitamente, já que é uma fonte recorrente de
citação incorreta). Aquele post prova resultados condicionais de
separação para $2^a-3^k$ via o método de Halász \cite{Halasz1977} e a
teoria inversa de Littlewood--Offord \cite{TaoVu2009,NguyenVu2011},
mas reconhece que o teorema de Baker já dá incondicionalmente uma
separação mais forte, tratando suas próprias proposições como
exercícios metodológicos. O problema clássico de ``Z-number'' de
Mahler \cite{Mahler1968} trata de $\xi(3/2)^n$, um objeto genuinamente
diferente (embora estruturalmente adjacente) do nosso (que trata de
$\times 2\bmod 3^s$ em profundidade finita, não $\times 3/2\bmod 1$
numa órbita infinita); a ancoragem honesta para nosso objeto é a
família de Erdős/Lagarias \cite{Erdos1979,Lagarias2009}.

\begin{definition}[Falha com margem $c$]
Para $c>0$, a Weak Covering Conjecture \emph{falha na escala $\ell$
com margem $c$} se, para $j=\lceil(1+c)\log_4 3\cdot\ell\rceil$, algum
$b\in(\mathbb{Z}/3^\ell\mathbb{Z})^\times$ está fora da imagem de
$R_{j-1,j}$. Defina $\gamma_c := (j+\ell)/\ell = 1+(1+c)\log_4 3$.
\end{definition}

\begin{definition}[Configuração diagonal]
Uma configuração de profundidade $s$ com inclinação $\gamma$,
tolerância $\delta$, e taxa de exceção $\eta$ é $\xi\in\ZZ$,
$3\nmid\xi$, tal que em $\ge(1-\eta)s$ das escalas $t\in\{1,\dots,s\}$
existe $a\in[\gamma t - t^{2/3}, \gamma t+t^{2/3}]$ com
$\|\xi\cdot 2^a/3^t\|<\delta$.
\end{definition}

\begin{lemma}[Trivialidade pré-wrap]
Para $\xi=1$, $\|2^a/3^s\|<\delta$ vale trivialmente para
$a\le\log_2 3\cdot s+\log_2\delta$ (o representante ainda não deu a
volta). Como $\gamma_c\ge 1{,}7925>\log_23\approx 1{,}585$ para todo
$c\ge 0$, a configuração forçada vive inteiramente no regime pós-wrap
não trivial, por folga estrutural vinda da contabilidade de entropia
$4^j$ vs.\ $3^\ell$.
\end{lemma}

\begin{definition}[Concentração espectral $SC(\varepsilon)$]
$\mu_{\ell,j}$ (a lei de $Y=\sum R\bmod 3^\ell$) satisfaz
$SC(\varepsilon)$ se existe $\xi\not\equiv 0\bmod 3^\ell$ com
$|\hat\mu_{\ell,j}(\xi)|\ge 3^{-\varepsilon\ell}$.
\end{definition}

\begin{lemma}[Redução condicional]\label{lem:B}
Para $c>0$ e $\varepsilon<\varepsilon_0(c)$ existem
$\delta(\varepsilon,c)<1/2$, $\eta(\varepsilon,c)<1$ tais que: falha
na escala $(\ell,c)$ junto com $SC(\varepsilon)$ implica existência de
uma configuração diagonal $(\gamma_c,\delta,\eta)$ de profundidade
$s\ge(1-O(\varepsilon))\ell$. Contrapositiva: ausência de configurações
diagonais de inclinação $\gamma_c$ em profundidade grande implica que
qualquer falha da Weak Covering Conjecture, se existir, é
espectralmente \emph{difusa} ($SC(\varepsilon)$ falha para todo
$\varepsilon$).
\end{lemma}

A prova (um argumento de seis etapas: cota inferior de inversão de
Fourier sobre $\sum|\hat\mu(\xi)|$, transferência para um modelo
i.i.d.\ inclinado no parâmetro $p_c=1/\gamma_c$, extração de
paralelepípedos alternáveis via o dispositivo de produto de Riesz de
\cite{Tao2011}, o teorema de Halász, uma descida de escala de rotina,
e um upgrade de uma configuração diagonal para um segmento contínuo)
está completa exceto pela última etapa, que exige estabilidade da
frequência relevante $\xi$ sob $j'\mapsto j'+1$ entre escalas ---
plausível mas não estabelecida, e declarada como lacuna aberta em vez
de dissimulada.

\subsection{Proposição C: uma parede exata de constantes}

\begin{proposition}[Déficit de Halász]\label{prop:halasz-deficit}
Incondicionalmente, o teorema de Halász dá apenas $\tau\ge 3^{-\ell}$,
não $SC(\varepsilon)$, produzindo um orçamento
$\sum\|\cdot\|^2\le 0{,}549\ell$ contra um teto necessário de
$0{,}25\ell$ ($\delta=1/2$, $d=\ell$ geradores): um déficit de fator
$\ge 2{,}2$, subindo a $\approx 7$ numa densidade realista de geradores
$\rho\approx 0{,}3$. Na aplicação original de Tao de 2011 isso fecha
porque $\log q \ll d$ (razão ilimitada); no regime da Weak Covering
Conjecture, $\log_2 q/d \ge 1{,}6$ --- a técnica não apenas degrada,
inverte de sinal.
\end{proposition}

\begin{theorem}[Identidade de Jensen e o déficit exato de orçamento de
Fourier]\label{thm:jensen}
Seja $Z=\sum_{g\ge1} w_g\, e(U_g)$, $U_g$ i.i.d.\ uniformes, $w_g =
p(1-p)^{g-1}$ ($\sum_g w_g=1$). Então, exatamente,
\begin{equation}\label{eq:jensen}
  \Lambda := \mathbb{E}\bigl[\log(1/|Z|)\bigr] = \log(1/p).
\end{equation}
\end{theorem}

\begin{proof}
Pela autossimilaridade sem memória dos pesos geométricos
($w_g/(1-w_1)=p(1-p)^{g-2}$ para $g\ge2$, a mesma sequência geométrica
reindexada), $Z \overset{d}{=} p\, e(U_1) + (1-p)\, Z'$, com $Z'$ uma
cópia independente de $Z$; como $|Z'|\le 1$ sempre (uma combinação
convexa de termos unimodulares), e $p/(1-p)\ge 1$ para $p\ge1/2$,
temos $|Z'|\le p/(1-p)$ sempre que $p\ge1/2$. Escrevendo $c:=(1-p)Z'/p$,
logo $|c|\le1$, fatore $p\,e(U_1)+(1-p)Z' = p\,e(U_1)\bigl(1+c\,e(-U_1)\bigr)$
e tome $\log|\cdot|$: a função $z\mapsto\log|1+cz|$ é (a parte real de)
uma função holomorfa de $z$ sem zero dentro do disco $|z|<1/|c|\ni$ o
círculo unitário (já que $|c|\le1$), então pela propriedade de valor
médio (fórmula de Jensen sem zeros envolvidos) sua média sobre
$z=e(U_1)$, $U_1$ uniforme, vale seu valor em $z=0$, ou seja $0$. Logo
$\mathbb{E}_{U_1}\log|p\,e(U_1)+(1-p)Z'| = \log p$ para toda
realização de $Z'$, e tomando a média sobre $Z'$ obtemos
$\mathbb{E}[\log|Z|]=\log p$, i.e.\ \eqref{eq:jensen}.
\end{proof}

Aplicando o Teorema~\ref{thm:jensen} com $p=p_c=1/\gamma_c$ (o modelo
inclinado que casa com a inclinação-limiar da Weak Covering
Conjecture) obtemos $\Lambda = \log\gamma_c \approx 0{,}5834+O(c)$,
confirmado por Monte Carlo ($8\times10^6$ amostras: $0{,}58344\pm0{,}0005$),
contra o valor $\log 3\approx 1{,}0986$ que seria necessário para
excluir um buraco difuso por um argumento anelado de orçamento de
Fourier $\ell^1$: um fator de déficit de $\approx 1{,}88$. O critério
anelado só inverteria em $\gamma=3$ ($j\ge 2\ell$), bem acima do regime
da conjectura ($j^\ast(\ell)/\ell\approx1{,}2$, Teorema~\ref{thm:wcc}).

\begin{theorem}[Proposição C: autoconsistência de um buraco
difuso]\label{thm:propC}
Uma falha espectralmente difusa da Weak Covering Conjecture satisfaz o
orçamento anelado de Fourier $\ell^1$ com folga $\approx 1{,}88$;
nenhum argumento de contagem de Littlewood--Offord/produto de Riesz do
tipo usado em \cite{Tao2011} pode excluí-la. Este ramo --- inacessível
a métodos $\ell^1$ pelo Teorema~\ref{thm:propC} --- é exatamente o
ramo identificado como o núcleo de $\beta=1$
(\S\ref{sec:triangulation}) e da condição $L^2$ refutada
(Teorema~\ref{thm:l2-refutation}): uma \emph{quarta e quinta}
articulação independente do mesmo ingrediente faltante.
\end{theorem}

Enfatizamos o enquadramento correto, consistente com as fragilidades
listadas acima: isto é uma redução do ramo espectralmente
\emph{concentrado} à família Erdős--Lagarias--Furstenberg, junto com
uma prova de que o ramo espectralmente \emph{difuso} é inacessível a
métodos $\ell^1$ de Littlewood--Offord --- não uma redução da Weak
Covering Conjecture ao problema clássico de Z-number, o que o
Teorema~\ref{thm:propC} e a discussão acima mostram que seria
impreciso.

\section{Contextualizando a barreira na literatura mais ampla}
\label{sec:context}

Uma busca literária dirigida, indo além da literatura específica de
Collatz para a teoria geral de decaimento de Fourier de medidas
autossimilares, situa a barreira com precisão e descarta vários atalhos
que pareciam plausíveis, cada um por uma razão específica e verificada.

\subsection{A medida de Syracuse é uma medida autossimilar $3$-ádica
rígida}

Desenrolando a recursão de renovação de \eqref{eq:syracuse-def} obtemos,
no limite,
\begin{equation}\label{eq:X-selfsimilar}
  X = \sum_{k\ge1} 3^{k-1}\, 2^{-S_k}, \qquad S_k:=a_1+\cdots+a_k,
\end{equation}
o ponto fixo do sistema de funções iteradas afim contável
$\phi_a(x) = 2^{-a}(1+3x)$ em $\mathbb{Z}_3$, ponderado por
$\Pr(a)=2^{-a}$, todos os mapas compartilhando a \emph{mesma} razão de
contração ultramétrica $|3\cdot2^{-a}|_3 = 1/3$. Esta é uma estrutura
autossimilar genuína (ainda que não-padrão, contável, com
sobreposições), então a questão de se a teoria geral de decaimento de
Fourier para medidas autossimilares \cite{BakerBanaji2026,Bremont,
LiSahlsten2022} se transfere é justa de se fazer, e a checamos
explicitamente.

\begin{proposition}[Rigidez de $\mathbb{Z}_3^\times$]\label{prop:rigidity}
Todo subgrupo fechado infinito de $(\mathbb{Z}_3,+)$ tem índice finito
(é igual a $3^k\mathbb{Z}_3$ para algum $k\ge 0$); consequentemente,
como $\mathbb{Z}_3^\times\cong\mathbb{Z}/2\times\mathbb{Z}_3$, todo
subgrupo fechado infinito de $\mathbb{Z}_3^\times$ tem índice finito.
Como $2$ é gerador topológico de $\mathbb{Z}_3^\times$ (raiz primitiva
módulo $3^n$ para todo $n$), as fases unitárias $2^{-a}$ se espalham ao
máximo, o oposto do fenômeno de ``órbita multiplicativa fina'' (uma
obstrução do tipo Pisot à propriedade de Rajchman) que impulsiona o
decaimento não uniforme de Fourier na reta real.
\end{proposition}

\begin{theorem}[A teoria geral não se transfere]\label{thm:no-transfer}
\begin{enumerate}[label=(\alph*)]
\item O mecanismo de pushforward de Baker--Banaji
  \cite{BakerBanaji2026} (fase estacionária arquimediana, exigindo um
  mapa $C^2$ com $f''\ne 0$ fora de um conjunto nulo) daria, no
  máximo, decaimento de Rajchman \emph{qualitativo} para o análogo
  considerado aqui --- estritamente mais fraco que o decaimento
  super-polinomial já garantido pela Proposição~\ref{prop:tao117}.
\item O critério de tipo Pisot de Brémont para medidas autossimilares
  não-Rajchman, transportado para $\mathbb{Z}_3$, trivializa pela
  Proposição~\ref{prop:rigidity}: não há espaço para uma ``órbita
  multiplicativa fina'' $3$-ádica do tipo que o critério exige.
\item A técnica de teoria de renovação de Li--Sahlsten
  \cite{LiSahlsten2022}, que precisa de \emph{duas} razões de
  contração distintas com uma relação diofantina entre seus
  logaritmos, também degenera: todo ramo de \eqref{eq:X-selfsimilar}
  tem a \emph{mesma} razão de contração $3$-ádica $1/3$, então o
  ``passeio de escalas'' relevante é determinístico ($S_n = n\log 3$,
  um passeio puramente lattice sem flutuação de overshoot), não o
  passeio aleatório não-lattice que a técnica é construída para
  explorar.
\end{enumerate}
\end{theorem}

Enfatizamos que o fenômeno geral de Baker--Banaji --- medidas de
Rajchman sem taxa de decaimento uniforme são genéricas, não
patológicas, na classe de medidas autossimilares/autoconformais ---
permanece um contexto justo: mostra que o tipo de falha exibido pelo
Teorema~\ref{thm:propC} tem precedente estrutural real na literatura
de referência, legitimando o enquadramento de que uma prova de
$\beta=1$/da Weak Covering Conjecture genuinamente precisa do insumo
aritmético específico deste problema, não de suavidade genérica. Não
deve, porém, ser sobre-interpretado como evidência \emph{a favor} do
cenário difuso em si: a construção de Baker--Banaji exige sintonia
fina de um parâmetro contínuo (uma construção do tipo Liouville, densa
na família mas excepcional em medida; parâmetros típicos dão
decaimento polinomial \cite{Solomyak1995,VarjuYu2022}), enquanto a
medida de Syracuse é um objeto aritmético rígido sem parâmetro
sintonizável, e a Proposição~\ref{prop:rigidity} mostra que o
mecanismo habilitador (aproximação diofantina fina) simplesmente não
existe do lado $3$-ádico --- o que corta contra o cenário difuso tanto
quanto a favor. Citamos o Teorema~\ref{thm:no-transfer} com essa
disanalogia declarada explicitamente.

\subsection{Uma sétima formulação independente, testada empiricamente}

Trabalho independente muito recente \cite{Chang2026} reduz a
conjectura de Collatz, por uma rota puramente combinatória sem
nenhuma análise de Fourier, a uma condição de balanço de resíduos ao
longo de uma subsequência esparsa da órbita verdadeira (não do
ensemble): escrevendo $X_t := \mathbb{1}[n_t\equiv1 \bmod 4]$ para a
órbita comprimida ímpar-a-ímpar e identificando ``tempos de
fim-de-rajada'' (o último índice de uma rajada maximal de $X_t=1$),
aquele trabalho prova um Teorema de Balanço do Mapa (o viés a nível de
mapa entre classes de comprimento de gap é exatamente $\pm1$, para
todo $K\ge5$) e reduz a conjectura, na classe de resíduo dominante, a
uma afirmação de que a fração de fins-de-rajada caindo num de dois
resíduos mod $32$ fica limitada perto de $1/2$ ao longo da órbita.

Testamos essa condição diretamente em órbitas reais. Um ensemble de
muitas órbitas moderadas ($2$--$5\times10^4$ raízes, $20$--$50$ bits)
mostra o desvio agregado estabilizando num platô de $1{,}2$--$1{,}9\%$,
o que por si só é consistente com a conjectura (só limitação é
afirmada, não convergência a zero); um teste mais fiel --- uma única
órbita muito longa (início aleatório de $16000$ bits, $38395$ passos
comprimidos, $4761$ fins-de-rajada relevantes) --- mostra o desvio
acompanhando de perto a curva de ruído estatístico $1/\sqrt{i}$, sem
tendência sistemática, terminando num desvio de $0{,}00053$. Cinco
outras órbitas de $500$ a $8000$ bits mostram o mesmo padrão.
Interpretamos o platô do ensemble como artefato de agregar muitas
órbitas curtas, finitas, de fase parecida, e o resultado da órbita
única longa como suporte empírico qualificado à conjectura de
\cite{Chang2026}, na escala testada.

Isto dá uma sétima articulação, tecnicamente não relacionada, do mesmo
ingrediente faltante: equidistribuição/balanço de resíduos ao longo da
órbita verdadeira, aparecendo independentemente em endogenia
(\S\ref{sec:endogeny}), na Weak Covering Conjecture (\S\ref{sec:wcc}),
em $\beta=1$ (\S\ref{sec:triangulation}), na condição $L^2$
(\S\ref{sec:lemmas}), na dicotomia espectral (\S\ref{sec:lemmas}), e
agora numa reformulação puramente combinatória.

\section{Discussão: por que uma barreira precisamente caracterizada é
uma contribuição genuína}
\label{sec:discussion}

Uma objeção natural ao material das \S\S\ref{sec:endogeny}--\ref{sec:context}
é que nenhum dos dois lemas candidatos da \S\ref{sec:lemmas} produziu
uma prova, então o paper poderia parecer documentar uma ausência em
vez de um resultado. Achamos que esta objeção, embora natural, não
sobrevive a um escrutínio do que de fato foi produzido.

Primeiro, o resultado de executar os dois lemas candidatos não foi
silêncio: foram dois resultados negativos/estruturais rigorosos e
independentes, cada um cravando \emph{precisamente} por que uma
abordagem natural falha --- a persistência exata em forma fechada de
correlação grosseira da Proposição~\ref{prop:exact-endogeny}; a
refutação computacional direta de uma condição suficiente específica
do Teorema~\ref{thm:l2-refutation}; a identidade exata do
Teorema~\ref{thm:jensen} quantificando, a um fator constante específico
($\approx1{,}88$), por que uma família inteira de técnicas de análise
de Fourier não consegue fechar a lacuna. Caracterizar precisamente uma
obstrução --- mostrar não apenas que uma abordagem falha, mas
\emph{exatamente} por quê, com uma margem quantificada --- é por si só
uma forma reconhecida e citável de contribuição matemática,
independente de a conjectura subjacente ser eventualmente resolvida.

Segundo, esses resultados não estão isolados: convergem, por sete
rotas tecnicamente não relacionadas (endogenia, a Weak Covering
Conjecture, $\beta=1$, a equação funcional de Wirsching de 2003, a
condição $L^2$, a dicotomia espectral da Proposição~C, e a
reformulação puramente combinatória de \cite{Chang2026}), no que
parece ser uma única afirmação aritmética subjacente: alguma forma de
equidistribuição ou quase-independência de dígitos $3$-ádicos ao longo
da dinâmica inteira verdadeira e determinística, não apenas ao longo
de um modelo i.i.d.\ idealizado. Uma barreira visível só de um ângulo
é candidata natural a ser artefato da técnica específica daquele
ângulo; uma barreira que recorre, essencialmente na mesma forma, em
sete vocabulários técnicos independentes (somas de caracteres,
transformações de suavização, uma equação funcional real, a identidade
de Jensen, uma classificação de medidas autossimilares, e contagem
elementar de resíduos) é evidência muito melhor de que a obstrução é
uma característica estrutural do problema em si.

Terceiro, este pacote teórico convive com um resultado
(Teorema~\ref{thm:kl}) cuja validade não depende de nada do
enquadramento acima: uma medição empírica direta, falsificável, e
estatisticamente calibrada, resolvendo uma disputa de quinze anos na
literatura. Consideramos esta a contribuição mais independente de
enquadramento do paper, e recomendamos que seja ponderada de acordo
por um leitor avaliando a confiabilidade geral do paper.

Apresentamos, portanto, a contribuição deste paper não como uma
tentativa de prova de qualquer parte da conjectura de Collatz, nem
como um relato de fracasso, mas como um mapa preciso de onde uma linha
de ataque específica, natural e tecnicamente rica --- o programa da
medida de Syracuse de Tao, generalizado para $qx+1$ --- esbarra numa
parede, junto com uma quantificação exata do tamanho e da forma dessa
parede a partir de seis ângulos independentes, e um resultado empírico
duro e incondicional obtido no caminho.

\section{Conclusão e direções em aberto}
\label{sec:conclusion}

Resumimos os problemas em aberto concretos que este paper deixa na
forma mais limpa que conseguimos dar a eles.

\begin{enumerate}[label=(A\arabic*)]
\item \textbf{Quase-independência de dígitos $3$-ádicos frescos entre
  subárvores irmãs} (Teorema~\ref{thm:barrier}): a afirmação
  aritmética cujo análogo idealizado (i.i.d.) é demonstrável via
  classificação de transformações de suavização e rigidez de
  Choquet--Deny.
\item \textbf{A Weak Covering Conjecture de Wirsching}
  (Conjectura~\ref{conj:wcc}), equivalentemente
  (Proposição~\ref{prop:beta-wcc}) o $\beta=1$ de Tao
  (\S\ref{sec:triangulation}), em aberto desde 1998 e 2020,
  respectivamente.
\item \textbf{As Conjecturas 1 e 2 de Wirsching} (\S\ref{sec:triangulation}),
  acima da Conjectura~3 testada numericamente, da mesma dificuldade de
  janela do teorema central do limite.
\item \textbf{Decaimento de Fourier exponencial, uniforme em $\xi$,
  para a medida de Syracuse em profundidade $\ell\asymp D$}
  (\S\ref{sec:regimes}, Regime~3), equivalente a problemas abertos
  reconhecidos em rigidez efetiva $\times2,\times3$ e decaimento de
  Fourier de medidas $3$-ádicas autossimilares.
\item \textbf{Exclusão de falhas espectralmente difusas da Weak
  Covering Conjecture} (Teorema~\ref{thm:propC}), mostradas aqui como
  inacessíveis a métodos $\ell^1$ de Littlewood--Offord, deixando
  técnica genuinamente nova como única rota.
\item \textbf{A conjectura de balanço de resíduos de Chang}
  \cite{Chang2026} (\S\ref{sec:context}), suportada empiricamente aqui
  numa única órbita muito longa, não provada.
\item \textbf{O índice de cauda do fator de escala do martingale}
  (Conjectura~\ref{conj:tail-index}), para todo $q\ge5$: diferente do
  caso $q=3$, isso não tem nem derivação rigorosa (a raiz relevante
  está sempre congelada, Proposição~\ref{prop:always-frozen}) nem apoio
  empírico sólido: a medição original, para $q=5$, já foi mostrada ser
  estatisticamente não-confirmatória, e um teste maior subsequente
  (\S\ref{sec:freezing}) é encorajador numa fração de cauda mas
  instável entre outras, portanto ainda não decisivo. Esta é,
  honestamente, a alegação com apoio mais fraco do paper, e a
  sinalizamos como tal em vez de deixá-la herdar credibilidade do caso
  $q=3$ só por compartilhar a fórmula.
\end{enumerate}

Qualquer progresso em qualquer um de (A1)--(A7) muito provavelmente se
transferiria, dadas as identidades estruturais estabelecidas neste
paper, para progresso em vários dos outros simultaneamente. Vemos esta
convergência como o achado central do paper.

\section*{Disponibilidade de código e dados}

Todo resultado computacional relatado neste paper --- os cálculos de
pressão e congelamento da \S\ref{sec:pressure}, o experimento de
controle do acoplamento residual da \S\ref{sec:calibration}, o portão
Kontorovich--Lagarias vs.\ Volkov da \S\ref{sec:kl}, o teste da Weak
Covering Conjecture da \S\ref{sec:wcc}, o teste da Conjectura~3 de
Wirsching (2003) e os cálculos de $L^2$/Jensen das
\S\ref{sec:triangulation}--\S\ref{sec:lemmas}, e o teste de balanço de
resíduos de Chang da \S\ref{sec:context} --- tem um script próprio,
independentemente executável, organizado por seção, em
\url{https://github.com/faculdade/collatz-endogeny}. Cada subpasta
daquele repositório tem seu próprio README dizendo o que verifica,
como rodar, e o resultado esperado.

\section*{Agradecimentos}

Esta pesquisa contou com o auxílio de sistemas de IA, utilizados para
apoio à revisão da literatura, análise matemática, experimentação
numérica e preparação do manuscrito. A referência a Halász
\cite{Halasz1977} não pôde ser verificada diretamente, pois o trabalho
original de 1971 não foi localizado; apenas sua errata de 1972 foi
encontrada. Por essa razão, a citação é mantida com a devida ressalva
quanto à impossibilidade de consulta à fonte primária. Seguimos, nesta
divulgação, o que acreditamos ser a melhor prática na pequena
literatura existente sobre contribuições assistidas por IA a este
problema específico.

\appendix

\section{Metodologia de validação numérica}
\label{app:validation}

Registramos, por completude, a disciplina de validação aplicada a cada
resultado computacional relatado acima, já que um cálculo recursivo
não validado é uma fonte conhecida de erro silencioso.

\begin{itemize}
\item \emph{Raízes da equação de pressão} (\S\ref{sec:pressure}):
  resolvidas por bisseção independentemente da derivação em forma
  fechada; conferidas contra inclinações de contagem empíricas em
  árvores reversas reais para $q=5,7$.
\item \emph{DP da Weak Covering Conjecture} (\S\ref{sec:wcc}):
  validada contra uma tabela de referência calculada
  independentemente para $\ell=1,\dots,7$ e um ponto conferido à mão
  ($\ell=2,j=2$).
\item \emph{Recursão de $K_\ell$} (Teorema~\ref{thm:l2-refutation}): a
  única checagem inicialmente disponível era o caso base hard-coded
  ($K_1=5/3$), que não exercita a maquinaria recursiva; validamos
  adicionalmente a distribuição completa $\mu_\ell$, bin a bin, contra
  amostragem direta de Monte Carlo da definição primária
  \eqref{eq:syracuse-def} em $\ell=3,4$, dentro do ruído de Monte
  Carlo.
\item \emph{Monte Carlo da identidade de Jensen} (Teorema~\ref{thm:jensen}):
  a identidade é exata e independente de simulação; a estimativa de
  Monte Carlo ($8\times10^6$ amostras) serve só como checagem de
  sanidade sobre o valor analítico, e bate com ele em quatro
  algarismos significativos.
\item \emph{Recursão de momentos de Wirsching 2003}
  (Teorema~\ref{thm:conjecture3}): validada por uma checagem
  independente de ponto fixo numa grade independente, concordando até
  a resolução da grade ($\approx 6$ dígitos), e confirmando a equação
  funcional~(7.12) de \cite{Wirsching2003} numericamente até
  $\ell=200$ usando uma derivada logarítmica analítica (não
  diferenciada numericamente).
\item \emph{Teste de balanço de resíduos de Chang} (\S\ref{sec:context}):
  validado confirmando que nenhum evento de fim-de-rajada jamais
  produziu um resíduo fora dos dois valores previstos ($9,25\bmod32$)
  nem no ensemble nem nos experimentos de órbita única, batendo
  exatamente com a proposição estrutural do paper-fonte.
\end{itemize}

\begin{thebibliography}{99}

\bibitem{Tao2022} T.\ Tao, \emph{Almost all orbits of the Collatz map
  attain almost bounded values}, Forum of Mathematics, Pi \textbf{10}
  (2022), e12. arXiv:1909.03562.

\bibitem{Tao2011} T.\ Tao, \emph{The Collatz conjecture, Littlewood--Offord
  theory, and powers of 2 and 3}, post de blog,
  \url{https://terrytao.wordpress.com/2011/08/25/}, 2011.

\bibitem{Tao2020} T.\ Tao, \emph{Equidistribution of Syracuse random
  variables and density of Collatz preimages}, post de blog,
  \url{https://terrytao.wordpress.com/2020/01/25/}, 2020.

\bibitem{Wirsching1998} G.\ J.\ Wirsching, \emph{The Dynamical System
  Generated by the $3n+1$ Function}, Lecture Notes in Mathematics
  \textbf{1681}, Springer, 1998.

\bibitem{Wirsching2003} G.\ J.\ Wirsching, \emph{On positive predecessor
  density in $3n+1$ dynamics}, Discrete and Continuous Dynamical
  Systems \textbf{9}(3) (2003), 771--787.

\bibitem{KontorovichLagarias2009} A.\ V.\ Kontorovich and J.\ C.\
  Lagarias, \emph{Stochastic models for the $3x+1$ and $5x+1$
  problems}, in: \emph{The Ultimate Challenge: the $3x+1$ Problem}
  (J.\ C.\ Lagarias, ed.), American Mathematical Society, 2010.

\bibitem{AldousBandyopadhyay2005} D.\ Aldous and A.\ Bandyopadhyay,
  \emph{A survey of max-type recursive distributional equations},
  Annals of Applied Probability \textbf{15}(2) (2005), 1047--1110.

\bibitem{GoncalvesGreenfeldMadrid2022} F.\ Gon\c{c}alves, R.\
  Greenfeld, and J.\ Madrid, \emph{Generalized Collatz maps with almost
  bounded orbits}, arXiv:2111.06170, 2022.

\bibitem{DurrettLiggett1983} R.\ Durrett and T.\ M.\ Liggett,
  \emph{Fixed points of the smoothing transformation}, Zeitschrift für
  Wahrscheinlichkeitstheorie und verwandte Gebiete \textbf{64}(3)
  (1983), 275--301.

\bibitem{AlsmeyerBigginsMeiners2012} G.\ Alsmeyer, J.\ D.\ Biggins, and
  M.\ Meiners, \emph{The functional equation of the smoothing
  transform}, Annals of Probability \textbf{40}(5) (2012), 2069--2105.

\bibitem{KoleskoMentemeier2015} K.\ Kolesko and S.\ Mentemeier,
  \emph{Fixed points of the multivariate smoothing transform: the
  critical case}, arXiv:1409.7220.

\bibitem{Goldie1991} C.\ M.\ Goldie, \emph{Implicit renewal theory and
  tails of solutions of random equations}, Annals of Applied
  Probability \textbf{1}(1) (1991), 126--166.

\bibitem{BuraczewskiDamekMikosch2016} D.\ Buraczewski, E.\ Damek, and
  T.\ Mikosch, \emph{Stochastic Models with Power-Law Tails: The
  Equation $X=AX+B$}, Springer, 2016.

\bibitem{KrasikovLagarias2003} I.\ Krasikov and J.\ C.\ Lagarias,
  \emph{Bounds for the $3x+1$ problem using difference inequalities},
  Acta Arithmetica \textbf{109} (2003), 237--258.

\bibitem{Biggins1992} J.\ D.\ Biggins, \emph{Uniform convergence of
  martingales in the branching random walk}, Annals of Probability
  \textbf{20}(1) (1992), 137--151.

\bibitem{ApplegateLagarias1995I} D.\ Applegate and J.\ C.\ Lagarias,
  \emph{Density bounds for the $3x+1$ problem I: tree-search method},
  Mathematics of Computation \textbf{64} (1995), 411--426.

\bibitem{ApplegateLagarias1995II} D.\ Applegate and J.\ C.\ Lagarias,
  \emph{Density bounds for the $3x+1$ problem II: Krasikov
  inequalities}, Mathematics of Computation \textbf{64} (1995),
  427--438.

\bibitem{Guivarch1990} Y.\ Guivarc'h, \emph{Sur une extension de la
  notion de loi semi-stable}, Annales de l'Institut Henri Poincaré
  Probabilités et Statistiques \textbf{26} (1990), 261--285.

\bibitem{Liu2000} Q.\ Liu, \emph{On generalized multiplicative
  cascades}, Stochastic Processes and their Applications
  \textbf{86}(2) (2000), 263--286.

\bibitem{Mentemeier2016} S.\ Mentemeier, \emph{The fixed points of the
  multivariate smoothing transform}, Probability Theory and Related
  Fields \textbf{164}(1--2) (2016), 401--458.

\bibitem{BakerBanaji2026} S.\ Baker and A.\ Banaji, \emph{Self-similar
  and self-conformal measures with slow Fourier decay},
  arXiv:2602.05593, 2026.

\bibitem{Bremont} J.\ Br\'emont, \emph{Self-similar measures and the
  Rajchman property}, Annales Henri Lebesgue \textbf{4} (2021),
  973--1004. arXiv:1910.03463.

\bibitem{LiSahlsten2022} J.\ Li and T.\ Sahlsten, \emph{Trigonometric
  series and self-similar sets}, Journal of the European Mathematical
  Society \textbf{24}(1) (2022), 341--368.

\bibitem{Solomyak1995} B.\ Solomyak, \emph{On the random series
  $\sum\pm\lambda^n$ (an Erd\H{o}s problem)}, Annals of Mathematics
  \textbf{142}(3) (1995), 611--625.

\bibitem{VarjuYu2022} P.\ Varj\'u and H.\ Yu, \emph{Fourier decay of
  self-similar measures and self-similar sets of uniqueness}, Analysis
  \& PDE \textbf{15}(3) (2022), 843--858. arXiv:2004.09358.

\bibitem{Spiegelhofer2021} L.\ Spiegelhofer, \emph{Collisions of digit
  sums in bases 2 and 3}, Israel Journal of Mathematics \textbf{258}
  (2023), 475--502. arXiv:2105.11173.

\bibitem{Chang2026} E.\ Y.\ Chang, \emph{A structural reduction of the
  Collatz conjecture to one-bit orbit mixing}, arXiv:2603.25753, 2026.

\bibitem{Halasz1977} G.\ Hal\'asz, \emph{On the distribution of
  additive and the mean values of multiplicative arithmetic
  functions}, Studia Scientiarum Mathematicarum Hungarica \textbf{6}
  (1971), 211--233.

\bibitem{TaoVu2009} T.\ Tao and V.\ Vu, \emph{Inverse Littlewood--Offord
  theorems and the condition number of random discrete matrices},
  Annals of Mathematics \textbf{169} (2009), 595--632.

\bibitem{NguyenVu2011} H.\ Nguyen and V.\ Vu, \emph{Optimal inverse
  Littlewood--Offord theorems}, Advances in Mathematics \textbf{226}
  (2011), 5298--5319.

\bibitem{Mahler1968} K.\ Mahler, \emph{An unsolved problem on the powers
  of $3/2$}, Journal of the Australian Mathematical Society
  \textbf{8}(2) (1968), 313--321.

\bibitem{FlattoLagariasPollington1995} L.\ Flatto, J.\ C.\ Lagarias,
  and A.\ D.\ Pollington, \emph{On the range of fractional parts
  $\{\xi(p/q)^n\}$}, Acta Arithmetica \textbf{70}(2) (1995), 125--147.

\bibitem{Erdos1979} P.\ Erd\H{o}s, \emph{Some unconventional problems
  in number theory}, Mathematics Magazine \textbf{52}(2) (1979),
  67--70.

\bibitem{Lagarias2009} J.\ C.\ Lagarias, \emph{Ternary expansions of
  powers of 2}, Journal of the London Mathematical Society
  \textbf{79}(3) (2009), 562--588.

\end{thebibliography}

\end{document}
